\documentclass[pdflatex,sn-mathphys-num]{sn-jnl}

\usepackage{amsmath,amssymb,amsfonts}%
\usepackage{amsthm}%
\usepackage{hyperref}%
\usepackage{enumitem}%

%% as per the template requirement, theorem styles follow sn-jnl's presets
\theoremstyle{thmstyleone}%
\newtheorem{theorem}{Theorem}[section]
\newtheorem{proposition}[theorem]{Proposition}
\newtheorem{lemma}[theorem]{Lemma}
\newtheorem{corollary}[theorem]{Corollary}

\theoremstyle{thmstylethree}%
\newtheorem{definition}[theorem]{Definition}

\theoremstyle{thmstyletwo}%
\newtheorem{remark}[theorem]{Remark}
\newtheorem{example}[theorem]{Example}
\newtheorem{openproblem}{Open Problem}

\DeclareMathOperator{\Sym}{Sym}
\DeclareMathOperator{\tr}{tr}
\DeclareMathOperator{\Hess}{Hess}
\DeclareMathOperator{\image}{im}
\DeclareMathOperator{\coker}{coker}
\DeclareMathOperator{\Diff}{Diff}
\DeclareMathOperator{\GL}{GL}
\DeclareMathOperator{\OO}{O}
\DeclareMathOperator{\Met}{Met}
\DeclareMathOperator{\id}{id}
\DeclareMathOperator{\Ric}{Ric}
\DeclareMathOperator{\Rm}{Rm}
\DeclareMathOperator{\Hol}{Hol}
\DeclareMathOperator{\SO}{SO}
\DeclareMathOperator{\diam}{diam}
\DeclareMathOperator{\Hom}{Hom}
\DeclareMathOperator{\Vol}{Vol}
\DeclareMathOperator{\Crit}{Crit}

\begin{document}

\title[Seam Geometry]{Seam Geometry: A Jet--Theoretic Framework for the Scalar Generation of Geometric Structures}

\author*[1]{\fnm{Lars} \sur{R\"onnb\"ack}}\email{lars@uptochange.com}

\affil*[1]{\orgdiv{Department of Computer and Systems Sciences}, \orgname{Stockholm University},
  \orgaddress{\city{Stockholm}, \country{Sweden}}}

\abstract{%
We develop a jet--theoretic framework for the seam--rule paradigm, in
which geometric structures are generated from scalar fields by local
coordinate-free transformation rules.  Viewing a rule as a \emph{natural
differential operator} on jet bundles yields a classification of all
isotropic metric-generating rules of order $\le 2$ on a Riemannian
background: an explicit three-parameter family consisting of a
conformal term, a gradient-tensor term, and a Hessian term, with
coefficients depending smoothly on scalar jet invariants.

We then develop the deformation theory of such rules.  The
linearisation at a seam is a linear differential operator whose
Fredholm theory controls local inverse design and whose kernel
represents gauge freedom.  We introduce the degeneration locus, the
geometrically invariant set where the linearised rule drops rank, and
explain how it unifies singular loci appearing in disparate
applications.

Finally, we study a universality problem: when can scalar seam data
generate an open neighbourhood in the space of Riemannian metrics (up
to diffeomorphism)?  We prove obstructions for order-$2$ rules,
establish a no-go theorem for ellipticity on surfaces when a Hessian
term is present, and prove local universality on the torus.
}

\keywords{seam geometry, natural differential operators, jet bundles, Hessian metrics, Teichm\"uller theory, universality}

\pacs[MSC Classification]{53C21 (primary), 58A20, 53A30, 53C25, 49Q20, 53C24}

\maketitle

% ================================================================
\section{Introduction}
% ================================================================

A recurrent pattern in geometry is the construction of rich geometric
structure from minimal data.  The pattern appears in diverse settings:
conformal and Hessian metrics on manifolds and meshes, Fisher--Rao
metrics in information geometry, Brenier potentials in optimal
transport, and score fields in generative diffusion models.  In every
case the construction follows the same schema:

\begin{center}
\emph{scalar field} $\;\xrightarrow{\;\text{rule}\;}\;$
\emph{geometric structure}.
\end{center}

We call the scalar field the \emph{seam} and the transformation the
\emph{rule}.  The present paper asks: \emph{what, precisely, is a rule?}

The naive answer---``a map from a scalar field and a background to a
geometric structure''---leaves the concept too vague to classify,
compose, or invert.  Our thesis is that the correct formalisation uses
\emph{jet bundles} and the theory of \emph{natural differential
operators}: a rule is a bundle map from the jet prolongation of the
trivial line bundle (encoding the seam and its derivatives up to some
order~$k$) to a target geometric bundle (encoding the output), subject
to a naturality (equivariance) condition.  This viewpoint, rooted in
the work of Kol\'a\v{r}, Michor, and Slov\'ak~\cite{KMS1993} and the
Peetre theorem~\cite{Peetre1960}, brings three immediate payoffs:

\begin{enumerate}[label=(\roman*)]
\item \textbf{Classification.}  The naturality condition, combined with
  invariant theory for the orthogonal group, constrains the space of
  possible rules to a tractable, finite-dimensional family at each jet
  order.  At order~$\le 2$, the family has exactly three generators
  (Theorem~\ref{thm:classification}).
\item \textbf{Linearisation and inverse design.}  The first variation
  of a rule at a seam is a linear differential operator whose
  Fredholm theory governs the solvability and stability of the inverse
  metric-design problem (\S\ref{sec:linearisation}).
\item \textbf{Universality.}  The image of the seam map under a given
  rule is a (typically proper) subset of all geometric structures.
  Characterising when this subset is dense or surjective is a
  well-posed geometric problem (\S\ref{sec:universality}).
\end{enumerate}

\begin{remark}[Relation to Hessian geometry and K\"ahler geometry]
\label{rem:prior-art}
Two well-established traditions exemplify special cases of the
seam--rule paradigm.  In \emph{Hessian geometry} (Shima and
collaborators), the metric on an affine manifold is defined as
$g = \nabla^2 \varphi$ for a convex potential~$\varphi$---precisely
the Hessian rule with $\alpha = \beta = 0$, $\gamma = 1$.  In
K\"ahler geometry, the $dd^c$-construction produces a K\"ahler
metric from a strictly plurisubharmonic potential; in real
coordinates this is again the Hessian of a convex function.  The
present framework subsumes these as particular coefficient choices
and contributes three things beyond what was previously known:
(a)~a proof that the conformal, gradient-tensor, and Hessian terms
are the \emph{only} generators at order~$\le 2$
(Theorem~\ref{thm:classification});
(b)~a systematic deformation theory and degeneration locus that
apply uniformly across all rules;
(c)~a higher-order extension (\S\ref{sec:higher}) that classifies
new generators at each jet order.
\end{remark}

% ================================================================
\section{Seams and Their Jets}
\label{sec:seams}
% ================================================================

Throughout, $(M^n,h)$ denotes a smooth Riemannian manifold of
dimension $n \ge 2$, and $C^\infty(M)$ (or an appropriate Sobolev
space) is the space of scalar functions on $M$.

\begin{definition}[Seam]
\label{def:seam}
A \emph{seam} on $(M,h)$ is a scalar field
$s \in \mathcal{S}(M) \subseteq C^\infty(M,\mathbb{R})$,
where $\mathcal{S}(M)$ is an admissible function class determined by
the application (e.g.\ $C^k$, Morse, convex, real-analytic, or
discrete vertex-based).
\end{definition}

The key observation is that any local transformation rule can depend on
$s$ only through its derivatives up to some finite order.  This is the
content of the classical Peetre theorem~\cite{Peetre1960}: every
support-non-increasing linear operator between sections of vector
bundles over $M$ is a differential operator of finite order.  For
nonlinear operators, the relevant generalisation is the Slov\'ak
theorem~\cite{Slovak1988}, which ensures that smooth natural operators
are determined by a finite jet.

\begin{definition}[$k$-jet of a seam]
\label{def:jet}
Let $s \in C^\infty(M)$.  The \emph{$k$-jet of $s$ at $x \in M$},
denoted $j^k_x s$, is the equivalence class of $s$ modulo functions
whose Taylor expansion at $x$ agrees through order $k$.  In local
coordinates $(x^1, \ldots, x^n)$,
\[
j^k_x s
= \bigl(s(x),\; \partial_i s(x),\;
  \partial_i \partial_j s(x),\; \ldots,\;
  \partial_{i_1} \cdots \partial_{i_k} s(x)\bigr).
\]
The collection of all $k$-jets forms a smooth fibre bundle
$J^k(M,\mathbb{R}) \to M$, the \emph{$k$-th jet bundle} of the
trivial line bundle.
\end{definition}

The jet bundle carries a natural $\GL(n)$-action induced by coordinate
changes on $M$.  If we restrict to orthonormal frames (i.e.\ use the
Riemannian metric $h$ to reduce the structure group), the relevant
group becomes $\OO(n)$.

\begin{remark}[Jet decomposition]
\label{rem:jet-decomposition}
At order $k = 2$, the jet $j^2_x s$ decomposes under the $\OO(n)$-action
(evaluated in $h$-normal coordinates at~$x$) into three irreducible
components:
\begin{enumerate}[label=(\alph*)]
\item The \emph{value}: $s_0 := s(x) \in \mathbb{R}$.
\item The \emph{gradient}: $s_i := \nabla_i s(x) \in T^*_x M \cong
  \mathbb{R}^n$.
\item The \emph{Hessian}: $s_{ij} := \nabla_i \nabla_j s(x) \in
  \Sym^2(T^*_x M)$, which further splits into
  $\tfrac{1}{n}(\Delta_h s)\,h_{ij}$ (trace) and the trace-free part.
\end{enumerate}
Understanding the $\OO(n)$-representation theory of this decomposition
is the key to classifying rules.
\end{remark}

% ================================================================
\section{Rules as Natural Differential Operators}
\label{sec:rules}
% ================================================================

\begin{definition}[Rule of order $k$]
\label{def:rule}
Let $\mathcal{G}(M)$ denote a natural fibre bundle over $M$ (the
\emph{geometric target}: e.g., $\Sym^2(T^*M)$ for Riemannian metrics,
$\Sym^2(T^*_\mathbb{C} M)$ for quadratic differentials, or
$\mathbb{R}^E$ for weighted graphs).  A \emph{rule of order $k$} is a
smooth bundle map
\begin{equation}
\label{eq:rule}
\Phi \colon J^k(M,\mathbb{R}) \times \mathcal{B}(M)
\longrightarrow \mathcal{G}(M)
\end{equation}
satisfying:
\begin{enumerate}[label=(R\arabic*)]
\item \label{item:locality}
  \textbf{Locality.}
  $\Phi$ depends fibrewise on
  $(j^k_x s, b_x)$ where $b_x$ is the fibre of the background
  structure $\mathcal{B}(M)$ at $x$.  That is, the output at $x$
  depends only on the $k$-jet of $s$ and the local background data
  at~$x$.
\item \label{item:naturality}
  \textbf{Naturality.}
  For every diffeomorphism $\varphi \colon M \to M$ preserving
  $\mathcal{B}$,
  \[
  \Phi(j^k(\varphi^* s), \varphi^* b)
  = \varphi^* \Phi(j^k s, b).
  \]
  If $\mathcal{B}$ includes a Riemannian metric $h$, naturality
  reduces to $\OO(n)$-equivariance at each fibre.
\item \label{item:regularity}
  \textbf{Regularity.}
  $\Phi$ is smooth (or at least continuous) as a map of
  fibre variables.
\end{enumerate}
\end{definition}

\begin{remark}[The three axioms explained]
Locality says that a rule ``looks at'' the seam through a finite-order
Taylor lens---it cannot see global information.
Naturality says the rule commutes with coordinate changes---it has no
preferred coordinate system.
Regularity prevents pathological constructions.
The first two axioms are precisely the sheaf-theoretic locality and
functoriality conditions familiar from discrete differential geometry,
now given their standard differential-geometric names.
\end{remark}

\begin{remark}[Naturality, diffeomorphisms, and cousin geometries]
\label{rem:cousin}
Naturality guarantees that if $\varphi\colon M \to M$ is a
diffeomorphism preserving the background and $s' = \varphi^* s$,
then $\Phi(s';\,\varphi^* h) = \varphi^*\Phi(s;\,h)$.
For the conformal rule $\Phi = e^{2s}h$, this means that
diffeomorphically related seams on the same background produce
isometric geometries whenever $\varphi$ is an isometry of~$h$.
However, for the Hessian rule $\Phi = \nabla^2 s$, a conformal
diffeomorphism $\varphi$ that is \emph{not} an isometry of~$h$
produces a genuinely different Hessian metric:
$\nabla^2(\varphi^* s) \neq \varphi^*(\nabla^2 s)$ in general,
because the Hessian depends on the background connection.

Concretely, the Poincar\'e half-plane seam $s_H = -\log y$
and the Poincar\'e disk seam $s_D = \ln(2/(1-|z|^2))$ produce
isometric conformal-rule geometries (both give curvature~$-1$),
but their Hessian-rule geometries are distinct: $\nabla^2 s_H$
is diagonal with a zero eigenvalue (capturing the vertical
foliation), while $\nabla^2 s_D$ is fully populated.
Thus \emph{the choice of seam representative within a
diffeomorphism orbit matters for non-conformal rules}---different
representatives probe different aspects of the same abstract
geometry.  This is analogous to gauge freedom: the conformal
rule is ``gauge-invariant,'' while the Hessian and gradient
rules carry additional information tied to the coordinate
representation of the seam.
\end{remark}

\begin{remark}[Connection to the Kol\'a\v{r}--Michor--Slov\'ak
  framework]
Definition~\ref{def:rule} is a special case of a \emph{natural operator}
in the sense of~\cite{KMS1993}: specifically, a natural operator from
sections of the trivial line bundle (with background structure
$\mathcal{B}$) to sections of the geometric bundle $\mathcal{G}$.
The general theory guarantees that the space of such operators at each
order is finite-dimensional and determined by $\GL(n)$-equivariant
maps between the corresponding standard fibres~\cite{KMS1993}.
The contribution of Theorem~\ref{thm:classification} is not the
abstract existence of such a finite-dimensional space (which follows
from the general KMS theory), but the \emph{explicit identification}
of the three generators together with the proof that no others exist
at order~$\le 2$, the systematic deformation theory and degeneration
locus analysis that the explicit form enables, and the extension
to higher jet orders (\S\ref{sec:higher}).
\end{remark}

\begin{definition}[The symbol of a rule]
\label{def:symbol}
The \emph{principal symbol} of a rule $\Phi$ of order $k$ at a
seam~$s$ is the leading-order part of $\Phi$ viewed as a map from
$S^k T^*_x M$ (the degree-$k$ symmetric cotangent tensors,
representing the ``pure order-$k$'' part of the jet) to the fibre of
$\mathcal{G}$ at $x$.  For a metric-valued rule, this is a map
\begin{equation}
\sigma_k(\Phi, s, x) \colon S^k T^*_x M
\longrightarrow \Sym^2(T^*_x M).
\end{equation}
The symbol governs the ellipticity properties of the linearised rule
and is the cornerstone of the inverse-design analysis
in~\S\ref{sec:linearisation}.
\end{definition}

% ================================================================
\section{Classification of Metric-Generating Rules}
\label{sec:classification}
% ================================================================

We now state the central structural result: a classification of all
isotropic natural differential operators from scalar fields to
Riemannian metrics at order $\le 2$.  The proof uses the classical
theory of $\OO(n)$-equivariant maps between representations.

\subsection{The representation-theoretic setup}

At a point $x \in M$ (working in $h$-normal coordinates), the 2-jet
of a scalar field provides the data
\[
(s_0,\; v,\; A)
\;\in\;
\mathbb{R} \times \mathbb{R}^n \times \Sym^2(\mathbb{R}^n),
\]
where $s_0 = s(x)$, $v_i = \nabla_i s(x)$, and
$A_{ij} = \nabla_i \nabla_j s(x)$.  The group $\OO(n)$ acts trivially
on $s_0$, by the standard representation on $v$, and by conjugation on
$A$.

The target is $\Sym^2(\mathbb{R}^n)$ (symmetric bilinear forms =
metrics).

By the \emph{first fundamental theorem of invariant theory} for
$\OO(n)$~\cite{Weyl1946,Procesi2007}, every polynomial
$\OO(n)$-equivariant map
\begin{equation}
\label{eq:equivariant-map}
F \colon \mathbb{R} \times \mathbb{R}^n \times \Sym^2(\mathbb{R}^n)
\longrightarrow \Sym^2(\mathbb{R}^n)
\end{equation}
can be expressed in terms of a finite set of basic equivariant maps
(the ``copies of $\Sym^2(\mathbb{R}^n)$ appearing in the tensor
product decomposition'') with polynomial-invariant coefficients.

\subsection{The three generators}

\begin{proposition}[Equivariant building blocks]
\label{prop:building-blocks}
The space of $\OO(n)$-equivariant \textbf{linear} maps
from $\mathbb{R} \oplus \mathbb{R}^n \oplus \Sym^2(\mathbb{R}^n)$ to
$\Sym^2(\mathbb{R}^n)$ is spanned by three generators:
\begin{enumerate}[label=(\roman*)]
\item $\Pi_h \colon (s_0, v, A) \mapsto s_0 \cdot \delta_{ij}$
  \quad\emph{(scalar-times-metric)},
\item $\Pi_\otimes \colon (s_0, v, A) \mapsto v_i v_j$
  \quad\emph{(gradient tensor)}---this is quadratic in $v$, hence
  contributes at the level of the ``squared 1-jet'',
\item $\Pi_{\Hess} \colon (s_0, v, A) \mapsto A_{ij}$
  \quad\emph{(Hessian)}.
\end{enumerate}
\end{proposition}

\begin{proof}[Proof sketch]
By Schur's lemma, an $\OO(n)$-equivariant linear map from an
irreducible representation $V$ to $\Sym^2(\mathbb{R}^n)$ is nonzero
only if $V$ contains $\Sym^2(\mathbb{R}^n)$ as a summand.
Now: $\mathbb{R}$ maps to $\Sym^2(\mathbb{R}^n)$ via the unique
copy of $\mathbb{R}$ in $\Sym^2(\mathbb{R}^n)$ (the trace component:
$\delta_{ij}$).  The symmetric square $S^2(\mathbb{R}^n)$ of the
standard representation decomposes as
$\mathbb{R} \oplus \Sym^2_0(\mathbb{R}^n)$, and the full
$\Sym^2(\mathbb{R}^n)$ matches this, so
$v_i v_j$ provides a quadratic equivariant map.
Finally, $\Sym^2(\mathbb{R}^n)$ maps to itself by the identity, giving
$A_{ij}$.
Higher-order terms (such as $(A^2)_{ij} = A_{ik} A_{kj}$, or
$v_i (Av)_j + v_j (Av)_i$) appear when we allow polynomial rather
than linear dependence on the 2-jet.
\end{proof}

\begin{theorem}[Universal Decomposition of order-$\le 2$ isotropic
  metric rules]
\label{thm:classification}
Let $(M^n, h)$ be a Riemannian manifold with $n \ge 2$.  Every smooth
isotropic natural differential operator
$\Phi \colon J^2(M, \mathbb{R}) \to \Sym^2_+(T^*M)$ that depends at
most linearly on the Hessian $\nabla^2 s$ takes the form
\begin{equation}
\label{eq:universal}
\boxed{
\Phi(j^2 s;\, h)
= \alpha\, h
  + \beta\, ds \otimes ds
  + \gamma\, \nabla^2 s,
}
\end{equation}
where $\alpha, \beta, \gamma$ are smooth real-valued functions of the
(finite list of) $\OO(n)$-invariant scalars of the 2-jet:
\begin{equation}
\label{eq:invariants}
\mathcal{I}
= \bigl(s,\;\; |\nabla s|_h^2,\;\;
  \Delta_h s,\;\;
  |\nabla^2 s|_h^2,\;\;
  \langle \nabla s, \nabla^2 s \cdot \nabla s \rangle_h
  \bigr).
\end{equation}
The condition $\Phi > 0$ (positive-definite output) imposes algebraic
constraints on $(\alpha, \beta, \gamma)$ at each point.

When we allow nonlinear dependence on the Hessian, additional
equivariant terms appear, namely
$(\nabla^2 s)^2 := (\nabla^2 s)_{ik}\,(\nabla^2 s)_{kj}$ and
$ds \odot (\nabla^2 s \cdot ds)$; the full ring of equivariants is
generated by these together with~\eqref{eq:universal}.
\end{theorem}

\begin{proof}
\emph{Step 1 (Reduction to fibre).}
By the Peetre--Slov\'ak theorem~\cite{Peetre1960,Slovak1988},
a smooth natural operator of order $\le 2$ factors through
$J^2(M,\mathbb{R})$.  By naturality (axiom~\ref{item:naturality}),
its value at a point depends only on the 2-jet of~$s$ at that point
and is $\OO(n)$-equivariant when expressed in $h$-normal coordinates.
The problem therefore reduces to classifying $\OO(n)$-equivariant maps
\[
F \colon \underbrace{\mathbb{R}}_{s_0}
\times \underbrace{\mathbb{R}^n}_{v}
\times \underbrace{\Sym^2(\mathbb{R}^n)}_{A}
\longrightarrow \Sym^2(\mathbb{R}^n)
\]
that are smooth and at most linear in $A$.

\emph{Step 2 (Invariant-theoretic decomposition).}
We decompose each factor as an $\OO(n)$-representation and identify
the equivariant maps to the target $\Sym^2(\mathbb{R}^n)$.

\begin{itemize}
\item \textbf{From}~$\mathbb{R}$ (trivial representation).
  $\Hom_{\OO(n)}(\mathbb{R}, \Sym^2(\mathbb{R}^n))
  \cong (\Sym^2(\mathbb{R}^n))^{\OO(n)}$.
  The $\OO(n)$-fixed subspace of $\Sym^2(\mathbb{R}^n)$ is
  one-dimensional, spanned by $\delta_{ij}$ (the identity matrix).
  This produces $s_0 \mapsto s_0\,\delta_{ij}$, i.e.\ the
  conformal generator~$\Pi_h$.

\item \textbf{From}~$\mathbb{R}^n$ (standard representation).
  There is no \emph{linear} $\OO(n)$-equivariant map
  $\mathbb{R}^n \to \Sym^2(\mathbb{R}^n)$, since
  $\Hom_{\OO(n)}(\mathbb{R}^n, \Sym^2(\mathbb{R}^n)) = 0$
  (the standard representation does not appear in
  $\Sym^2(\mathbb{R}^n) \cong \mathbb{R} \oplus
  \Sym^2_0(\mathbb{R}^n)$).  However, the \emph{quadratic}
  equivariant $v \mapsto v \otimes v$ maps
  $S^2(\mathbb{R}^n) \to \Sym^2(\mathbb{R}^n)$ by the identity on
  $S^2(\mathbb{R}^n) \cong \Sym^2(\mathbb{R}^n)$.
  By the first fundamental theorem for
  $\OO(n)$~\cite{Weyl1946,Procesi2007}, this is the unique
  quadratic equivariant (up to coefficient), producing~$\Pi_\otimes$.

\item \textbf{From}~$\Sym^2(\mathbb{R}^n)$ (target copy).
  As an $\OO(n)$-module,
  $\Sym^2(\mathbb{R}^n) = \mathbb{R}\cdot\delta
  \oplus \Sym^2_0(\mathbb{R}^n)$
  where $\Sym^2_0$ denotes the traceless part, which is irreducible
  for $n \geq 2$.  By Schur's lemma,
  $\Hom_{\OO(n)}(\Sym^2(\mathbb{R}^n), \Sym^2(\mathbb{R}^n))
  \cong \mathbb{R} \oplus \mathbb{R}$,
  spanned by the identity $A \mapsto A$ and the trace map
  $A \mapsto (\tr A)\,\delta$.  The trace map
  contributes $\Delta s \cdot h$, which is already of the
  form $\alpha(\Delta s)\,h$ and thus absorbed into the
  conformal generator with $\alpha$ depending on $\Delta s$.
  The identity map gives the Hessian generator~$\Pi_{\Hess}$.
\end{itemize}

\emph{Step 3 (Coefficients).}
The first fundamental theorem~\cite{Weyl1946} tells us that
the scalar $\OO(n)$-invariants of the 2-jet $(s_0, v, A)$ are
generated by $s_0$, $|v|^2 = v_i v_i$, $\tr A = A_{ii}$,
$|A|^2 = A_{ij}A_{ij}$, and $\langle v, Av\rangle = v_i A_{ij} v_j$.
The coefficients $\alpha, \beta, \gamma$ may be arbitrary smooth
functions of these invariants.  On $M$, these are
exactly~\eqref{eq:invariants}.

\emph{Step 4 (Completeness at the linear level).}
Every equivariant map from $\mathbb{R} \oplus \mathbb{R}^n \oplus
\Sym^2(\mathbb{R}^n)$ to $\Sym^2(\mathbb{R}^n)$ that is linear
in $A$ and polynomial of degree $\leq 2$ in $v$ is a linear
combination (with invariant coefficients) of the three generators
above.  This follows from the exhaustive decomposition in Step~2:
there are no further copies of $\Sym^2(\mathbb{R}^n)$ in the
relevant tensor products.

\emph{Step 5 (Nonlinear extension).}
Allowing polynomial dependence on $A$ beyond degree~1 introduces
higher equivariants.  The symmetric square $S^2(\Sym^2(\mathbb{R}^n))$
contains $\Sym^2(\mathbb{R}^n)$ with multiplicity~2, yielding
$(A^2)_{ij} = A_{ik}A_{kj}$ and $(\tr A)\,A_{ij}$ (the latter
already captured by an invariant-dependent coefficient on $\Pi_{\Hess}$).
The mixed tensor product $\mathbb{R}^n \otimes \Sym^2(\mathbb{R}^n)
\otimes \mathbb{R}^n$ contributes the symmetrised product
$v_i(Av)_j + v_j(Av)_i = ds \odot (\nabla^2 s \cdot ds)$.
These exhaust the generators through degree~2 in~$A$.
\end{proof}

\begin{corollary}[Forced factorisation through invariants]
\label{cor:forced-arch}
Under the hypotheses of Theorem~\ref{thm:classification}, the value of
$\Phi(j^2s;h)$ at a point $x\in M$ depends on the 2-jet of $s$ at $x$
only through the five invariant scalars $\mathcal{I}(x)$ in
\eqref{eq:invariants} and the three basis tensors
$h(x)$, $ds\otimes ds(x)$, and $\nabla^2 s(x)$.

Equivalently, any parametric family of such rules (e.g. a
``metric network'' constrained to represent a smooth isotropic natural
operator of order $\le 2$ that is at most linear in $\nabla^2s$) can be
written in the canonical form
\[
\Phi_\theta(j^2 s;h)
= \alpha_\theta(\mathcal{I})\,h
  + \beta_\theta(\mathcal{I})\,ds\otimes ds
  + \gamma_\theta(\mathcal{I})\,\nabla^2 s,
\]
for some scalar-valued coefficient functions
$(\alpha_\theta,\beta_\theta,\gamma_\theta)\colon \mathbb{R}^5\to\mathbb{R}^3$.
\end{corollary}

\begin{proof}
This is exactly the content of the decomposition
\eqref{eq:universal} together with the description of admissible
coefficient dependence in~\eqref{eq:invariants}.
\end{proof}

\begin{remark}[Relation to scalar-tensor theories in physics]
\label{rem:horndeski}
The decomposition in Theorem~\ref{thm:classification} has a parallel
in theoretical physics: scalar-tensor gravity theories (Horndeski~\cite{Horndeski1974},
Galileon, and beyond-Horndeski Lagrangians) employ the same
$\OO(n)$-invariant building blocks $h$, $ds \otimes ds$, and
$\nabla^2 s$ to construct covariant actions from a scalar field
and a metric.  The classification there is driven by the requirement
of second-order field equations rather than naturality of the output,
but the underlying invariant-theory content is closely related.
The present jet-theoretic treatment makes the representation-theoretic
reason for the three generators fully explicit
(Steps~1--5 above) and extends naturally to higher jet orders
(\S\ref{sec:higher}).
\end{remark}

\subsection{Recovery of known rules}

Every concrete rule encountered in practice is a specialisation of
\eqref{eq:universal}:

\begin{example}[Conformal Rule]
\label{ex:conformal}
$\alpha(s) = e^{2s}$, \; $\beta = 0$, \; $\gamma = 0$.
Output: $\Phi = e^{2s} h$.
\end{example}

\begin{example}[Warped / Canalization Rule]
$\alpha(s) = \phi(s)^2$ for an admissible warping function $\phi$, \;
$\beta = \gamma = 0$.
Output: $\Phi = \phi(s)^2 h$.
\end{example}

\begin{example}[Hessian Rule]
$\alpha = 0$, \; $\beta = 0$, \; $\gamma = 1$.
Output: $\Phi = \nabla^2 s$.
\end{example}

\begin{proposition}[No globally Riemannian pure Hessian metrics on closed manifolds]
\label{prop:hess-no-go-closed}
Let $(M^n,h)$ be closed (compact without boundary) and let $s\in C^\infty(M)$.
Then the symmetric 2-tensor $\nabla_h^2 s$ cannot be positive-definite at
every point.  Equivalently, the pure Hessian rule $g=\nabla_h^2 s$ cannot
produce a globally Riemannian metric on a closed manifold.
\end{proposition}

\begin{proof}
If $\nabla_h^2 s$ were positive-definite everywhere, then its $h$-trace
$\Delta_h s=\tr_h(\nabla_h^2 s)$ would satisfy $\Delta_h s>0$ pointwise.
But on a closed manifold, integration by parts gives
$\int_M \Delta_h s\,dV_h=0$, a contradiction.
\end{proof}

\begin{example}[Gradient (Eikonal) Rule]
$\alpha = |\nabla s|_h^2$, \; $\beta = 0$, \; $\gamma = 0$.
Output: $\Phi = |\nabla s|^2 h$.
\end{example}

\begin{example}[Mixed conformal--Hessian rule (new)]
\label{ex:mixed}
$\alpha(s) = e^{2s}$, \; $\beta = 0$, \; $\gamma(\Delta s) =
-\tfrac{1}{n-2}\,e^{2s}$ (for $n \ge 3$).
Output:
$\Phi = e^{2s}\bigl(h - \tfrac{1}{n-2}\nabla^2 s\bigr)$.
Note that the Schouten tensor of the conformal metric $g = e^{2s}h$
in dimension $n$ involves exactly this combination of $h$ and
$\nabla^2 s$ (with the Ricci correction absorbed through $s$);
thus the Schouten tensor of a conformal metric is itself
a seam-generated object via a mixed rule.
\end{example}

\begin{corollary}[Schouten tensor change is a universal seam rule]
\label{cor:schouten}
Let $n\ge 3$ and let $g=e^{2s}h$.
Writing $P_g := \frac{1}{n-2}\bigl(\Ric_g - \frac{R_g}{2(n-1)}g\bigr)$
for the Schouten tensor, the conformal change law may be written as
\[
P_h - P_g
= \nabla^2 s - ds\otimes ds + \tfrac12|\nabla s|_h^2\,h.
\]
In particular, $P_h-P_g$ is the output of a universal three-term rule
applied to $s$ with coefficient triple
$(\alpha,\beta,\gamma)=(\tfrac12|\nabla s|_h^2,\,-1,\,1)$.
\end{corollary}

\begin{proof}
This is the standard conformal transformation formula for the Schouten
tensor, rewritten in the three-term basis.
\end{proof}

\begin{remark}[What the classification reveals]
\label{rem:what-classif}
The Universal Decomposition shows that the three ``basis directions''
$h$, $ds \otimes ds$, and $\nabla^2 s$ exhaust the linear-in-Hessian
possibilities.  In particular:
\begin{itemize}
\item \emph{Conformal rules use only the $h$-direction}---they
  cannot see the gradient tensor or the Hessian directly.
\item \emph{The gradient-tensor direction $ds \otimes ds$ has been
  largely unexplored} in the literature.  It produces
  rank-deficient (degenerate at critical points) metrics and is
  related to sub-Riemannian geometry and Carnot--Carath\'eodory
  structures.
\item \emph{Mixed rules} (nonzero $\alpha$ and $\gamma$
  simultaneously) capture curvature-corrected geometries and arise
  naturally in conformal geometry (Example~\ref{ex:mixed}).
\end{itemize}
\end{remark}

\begin{corollary}[Order-1 rules preserve the level-set foliation]
\label{cor:levelset-foliation}
Let $g=\alpha\,h+\beta\,ds\otimes ds$ be a seam metric (so $\gamma=0$)
with $\alpha>0$ and $\alpha+\beta|\nabla s|_h^2>0$.
Then the $g$-gradient of $s$ is a pointwise scalar multiple of the
$h$-gradient:
\[
\nabla^g s = \frac{1}{\alpha+\beta|\nabla s|_h^2}\,\nabla^h s.
\]
Consequently, the orthogonal distribution $\ker(ds)$ (the tangent
spaces to level sets of $s$) is the same for $h$ and for $g$, and the
gradient flow lines of $s$ are unchanged up to reparametrisation.
\end{corollary}

\begin{proof}
The inverse of a rank-1 update is explicit:
$g^{-1}=\alpha^{-1}h^{-1}-\frac{\beta}{\alpha(\alpha+\beta|\nabla s|_h^2)}\,\nabla^h s\otimes\nabla^h s$.
Applying this to $ds$ gives the stated formula for $\nabla^g s$.
\end{proof}

\begin{corollary}[Volume form for order-1 rules]
\label{cor:volume-order1}
Let $g=\alpha\,h+\beta\,ds\otimes ds$ with $\alpha>0$ and
$\alpha+\beta|\nabla s|_h^2>0$.  Then
\[
dV_g
= \alpha^{\frac{n-1}{2}}\,\sqrt{\alpha+\beta|\nabla s|_h^2}\;dV_h
= \alpha^{\frac n2}\,\sqrt{1+\frac{\beta}{\alpha}|\nabla s|_h^2}\;dV_h.
\]
\end{corollary}

\begin{proof}
In $h$-orthonormal coordinates, $g=\alpha I+\beta v v^\top$ with
$|v|^2=|\nabla s|_h^2$.  The matrix determinant lemma gives
$\det(\alpha I+\beta v v^\top)=\alpha^n\bigl(1+\frac{\beta}{\alpha}|v|^2\bigr)$.
Taking square roots and restoring $dV_h$ yields the claim.
\end{proof}

\begin{corollary}[Determinant and degeneracy for order-1 rules]
\label{cor:order1-degeneracy}
Let $g=\alpha\,h+\beta\,ds\otimes ds$ be an order-$1$ seam tensor.
Then at every point,
\[
\det(g)=\alpha^{n-1}\bigl(\alpha+\beta|\nabla s|_h^2\bigr)\,\det(h).
\]
In particular, $g$ is nondegenerate if and only if
$\alpha\neq 0$ and $\alpha+\beta|\nabla s|_h^2\neq 0$.
Moreover, if $\alpha>0$ then $g$ is Riemannian if and only if
$\alpha+\beta|\nabla s|_h^2>0$.
\end{corollary}

\begin{proof}
The determinant identity is exactly the computation in the proof of
Corollary~\ref{cor:volume-order1}, rewritten as
$\det(g)=\det(h)\,\det(\alpha I+\beta v v^\top)$ in an $h$-orthonormal frame.
The nondegeneracy and positivity conditions follow immediately.
\end{proof}

\begin{proposition}[Signature and generic degeneracy hypersurfaces for order-1 rules]
\label{prop:order1-signature}
Let $g=\alpha\,h+\beta\,ds\otimes ds$ with $\alpha>0$ and define the scalar
function $f:=\alpha+\beta|\nabla s|_h^2$.
\begin{enumerate}[label=\textup{(\alph*)}]
\item $g$ is Riemannian at $x$ if and only if $f(x)>0$.
\item If $f(x)<0$, then $g$ has Lorentzian signature $(n-1,1)$ at $x$.
\item If $f(x)=0$, then $g$ has rank $n-1$ at $x$, and its kernel is the
  line spanned by $\nabla^h s(x)$.
\item If $0$ is a regular value of $f$, then the degeneracy set
  $\Sigma:=f^{-1}(0)$ is a smooth embedded hypersurface.
\end{enumerate}
\end{proposition}

\begin{proof}
Fix $x\in M$.  If $\nabla s(x)=0$ then $g(x)=\alpha(x)h(x)$, so $f(x)=\alpha(x)>0$
and $g$ is Riemannian.
Assume $\nabla s(x)\neq 0$ and choose an $h$-orthonormal basis
$\{e_1,\ldots,e_n\}$ with $e_1=\nabla^h s/|\nabla s|_h$.  Then
$ds=|\nabla s|_h\,e^1$ and
\[
g
=\alpha\sum_{i=1}^n e^i\otimes e^i + \beta|\nabla s|_h^2\,e^1\otimes e^1,
\]
so the $g$-eigenvalues are $\alpha$ with multiplicity $n-1$ on
$\mathrm{span}\{e_2,\ldots,e_n\}$ and $\alpha+\beta|\nabla s|_h^2=f$ on $e_1$.
This proves (a)--(c).  Part~(d) is the regular value theorem.
\end{proof}

\begin{proposition}[Comparison for order-1 seam metrics]
\label{prop:order1-comparison}
Let $g_i=\alpha_i\,h+\beta_i\,ds\otimes ds$ ($i=1,2$) be order-$1$ seam
metrics with $\alpha_i>0$ and $f_i:=\alpha_i+\beta_i|\nabla s|_h^2>0$.
Assume pointwise that $\alpha_1\le \alpha_2$ and $f_1\le f_2$.  Then:
\begin{enumerate}[label=\textup{(\alph*)}]
\item $g_1\le g_2$ as quadratic forms, hence $d_{g_1}(x,y)\le d_{g_2}(x,y)$
  for all $x,y\in M$.
\item $dV_{g_1}\le dV_{g_2}$ pointwise, hence $\Vol_{g_1}(M)\le \Vol_{g_2}(M)$.
\item If $c$ is a regular value of $s$ and $\Sigma_c:=\{s=c\}$, then the
  induced $(n{-}1)$-volume forms satisfy
  $dA_{g_1|_{\Sigma_c}}\le dA_{g_2|_{\Sigma_c}}$ pointwise on $\Sigma_c$.
  In particular, $\Vol_{n-1}(\Sigma_c,g_1)\le \Vol_{n-1}(\Sigma_c,g_2)$.
\end{enumerate}
\end{proposition}

\begin{proof}
Fix $x\in M$.  If $\nabla s(x)=0$ then $g_i(x)=\alpha_i(x)h(x)$ and
$\alpha_1\le\alpha_2$ gives $g_1\le g_2$.  If $\nabla s(x)\neq 0$, choose an
$h$-orthonormal basis with $e_1=\nabla^h s/|\nabla s|_h$.  In this basis,
$g_i$ is diagonal with eigenvalue $f_i$ on $e_1$ and eigenvalue $\alpha_i$
on $\mathrm{span}\{e_2,\ldots,e_n\}$.  The inequalities
$\alpha_1\le\alpha_2$ and $f_1\le f_2$ imply $g_1\le g_2$ pointwise.

For (a), for any piecewise smooth curve $\gamma$ we have
$|\dot\gamma|_{g_1}\le |\dot\gamma|_{g_2}$ pointwise, hence
$\mathrm{Len}_{g_1}(\gamma)\le \mathrm{Len}_{g_2}(\gamma)$; taking the
infimum over curves yields $d_{g_1}\le d_{g_2}$.

For (b), the volume-form identity of Corollary~\ref{cor:volume-order1}
gives $dV_{g_i}=\alpha_i^{\frac{n-1}{2}}\sqrt{f_i}\,dV_h$, so the assumed
inequalities imply $dV_{g_1}\le dV_{g_2}$.

For (c), $g_1\le g_2$ implies $g_1|_{T\Sigma_c}\le g_2|_{T\Sigma_c}$ as
quadratic forms on $T\Sigma_c$.  Hence in any basis of $T\Sigma_c$,
$\det(g_1|_{T\Sigma_c})\le \det(g_2|_{T\Sigma_c})$, which is equivalent to
$dA_{g_1|_{\Sigma_c}}\le dA_{g_2|_{\Sigma_c}}$.
\end{proof}

\begin{proposition}[Geometric constraints by rule type]
\label{prop:holonomy}
Let $(M^n,h)$ be simply connected and compact, and let
$g = \Phi(s;h)$ be a seam metric.  Then:
\begin{enumerate}[label=\textup{(\alph*)}]
\item \textbf{Conformal rule}
  ($\alpha = e^{2s}$, $\beta = \gamma = 0$).
  The metric $g = e^{2s}h$ is conformally equivalent to~$h$,
  so the conformal class $[g] = [h]$ is preserved.
  In particular, any invariant of the underlying conformal
  structure (i.e.\ depending only on $[h]$) is shared by~$g$.
  Note, however, that conformal rescaling
  generically changes the Levi-Civita connection and can alter
  the holonomy group: if $\Hol(h)$ is a proper subgroup of
  $\SO(n)$ (e.g.\ arising from a Ricci-flat or K\"ahler
  structure), then $g = e^{2s}h$ typically has
  $\Hol(g) = \SO(n)$ unless $s$ satisfies additional
  compatibility conditions.
\end{enumerate}
\end{proposition}

\begin{proof}[Proof sketch]
Part~(a).  The Levi-Civita connections of $g = e^{2s}h$ and $h$
differ by an exact 1-form: $\nabla^g = \nabla^h + ds \otimes
\id + \id \otimes ds - h(\cdot,\cdot)\nabla s$.  This is
the standard conformal change formula.  Since $g$ and $h$
differ by a positive scalar factor, they share the same
conformal class, Weyl tensor (for $n \ge 4$), and conformal
invariants.  The connection itself, however, changes
non-trivially, so holonomy is not preserved in general.
\end{proof}

% ================================================================
\section{Worked Examples}
\label{sec:worked}
% ================================================================

We now develop four concrete computations that illustrate the
Universal Decomposition in action.  Each takes a specific seam and
rule, computes the resulting geometry explicitly, and highlights the
features that make the framework non-trivial.

% ----------------------------------------------------------------
\subsection{The round sphere via a conformal seam}
\label{sec:ex-sphere}
% ----------------------------------------------------------------

Let $M = \mathbb{R}^2$ with the flat metric $h = \delta$, and define
the seam
\[
s(x) = -\log(1 + |x|^2).
\]
Applying the conformal rule $(\alpha, \beta, \gamma) = (e^{2s}, 0, 0)$
gives
\[
g = e^{2s}\,\delta
  = \frac{1}{(1+|x|^2)^2}\,
    (dx^1 \otimes dx^1 + dx^2 \otimes dx^2).
\]
This is exactly the round metric on $S^2$ in stereographic
coordinates (up to a factor of~$4$ depending on the radius
convention).

\paragraph{Curvature.}
The Gauss curvature of a conformal metric $g = e^{2s}\delta$ on
$\mathbb{R}^2$ is
\[
K = -e^{-2s}\,\Delta s.
\]
A direct computation gives
$\Delta s = -\frac{4}{(1+|x|^2)^2}$, so
\[
K = -\,(1+|x|^2)^2 \cdot
    \Bigl(-\frac{4}{(1+|x|^2)^2}\Bigr) = 4.
\]
Thus the seam encodes a surface of constant positive curvature.
This is one of the simplest non-flat geometries, yet the seam that
produces it is just a radial logarithm.

\paragraph{Degeneration.}
Since $e^{2s} = (1+|x|^2)^{-2} > 0$ everywhere, the conformal rule
has no degeneration locus for this seam.  This is consistent with the
general principle: the conformal rule $\alpha(s) = e^{2s}$ never
degenerates, because $e^{2s}$ has no zeros.

\paragraph{Linearisation.}
A seam perturbation $\dot s$ at $s_0 = s$ produces the metric
variation $\dot g = 2\dot s\, g$, i.e.\ a pure conformal
deformation.  On a compact surface, the only $\dot s$ that leaves
$g$ unchanged is $\dot s = 0$: the kernel of the linearised
conformal rule is trivial, confirming that the seam faithfully
parametrises the geometry.

% ----------------------------------------------------------------
\subsection{Lorentzian geometry from a harmonic seam}
\label{sec:ex-hessian}
% ----------------------------------------------------------------

Let $M = \mathbb{R}^2$ with $h = \delta$.  Consider the seam
\[
s(x,y) = x^2 - y^2 = \operatorname{Re}(z^2),
\]
which is harmonic ($\Delta s = 0$).  The Hessian rule
$(\alpha,\beta,\gamma) = (0,0,1)$ gives
\[
g = \nabla^2 s
  = \begin{pmatrix} 2 & 0 \\ 0 & -2 \end{pmatrix}
    dx^i \otimes dx^j.
\]

\paragraph{Signature.}
The eigenvalues $(+2, -2)$ show that $g$ has \emph{Lorentzian}
(indefinite) signature.  The Hessian rule applied to a harmonic seam
on $\mathbb{R}^2$ always produces a Lorentzian metric: harmonicity
forces $\Delta s = \tr(\nabla^2 s) = 0$, which means the eigenvalues
sum to zero and therefore have opposite signs (away from $\nabla^2 s = 0$).

This is the mechanism by which the Hessian rule, applied to seams
derived from holomorphic functions (whose real parts are harmonic),
produces Lorentzian geometry.

\paragraph{Geodesics.}
The metric $g = 2\,dx^2 - 2\,dy^2$ is flat Minkowski space
(rescaled).  Null geodesics satisfy $dx^2 = dy^2$, i.e.\
$y = \pm x + c$: they are 45-degree lines in the $(x,y)$-plane.
These are precisely the level curves $\operatorname{Re}(z^2) =
\text{const}$ and $\operatorname{Im}(z^2) = \text{const}$, which
form a family of orthogonal hyperbolas.

\paragraph{Degeneration.}
The Hessian $\nabla^2 s$ is constant and nondegenerate, so the
degeneration locus is empty for this particular seam.  By contrast,
for the seam $s = \log|f(z)|$ with $f$ holomorphic, the Hessian
degenerates precisely at the zeros of $f$---an example that
illustrates why the degeneration locus is the geometric locus of
interest.

% ----------------------------------------------------------------
\subsection{Score geometry of a Gaussian}
\label{sec:ex-gradient}
% ----------------------------------------------------------------

Let $M = \mathbb{R}^n$ with $h = \delta$, and let $p(x) \propto
\exp(-\|x\|^2/(2\sigma^2))$ be a centred Gaussian density.  The
log-density seam is
\[
s(x) = \log p(x) = -\frac{\|x\|^2}{2\sigma^2} + C
\]
for a normalisation constant $C$.  The gradient is
$\nabla s = -x/\sigma^2$, so $|\nabla s|^2 = \|x\|^2/\sigma^4$.
The gradient rule $(\alpha,\beta,\gamma) = (|\nabla s|^2, 0, 0)$
produces
\[
g = \frac{\|x\|^2}{\sigma^4}\,\delta.
\]

\paragraph{Geometric meaning.}
This is a conformal metric on $\mathbb{R}^n$ whose conformal factor
$\|x\|^2/\sigma^4$ vanishes at the origin and grows quadratically.
Distances near the origin are compressed to zero; distances far from
the origin are stretched.  Intuitively, the geometry ``funnels'' all
paths toward the mode of the Gaussian.

\paragraph{Degeneration at the mode.}
The degeneration locus is $\{x = 0\}$: the single point where
$\nabla s = 0$, i.e.\ the mode of the distribution.  At this point
the metric $g$ vanishes, and the Riemannian distance from any $x
\ne 0$ to the origin is
\[
d_g(x, 0) = \int_0^1 \frac{\|tx\|}{\sigma^2}\,\|x\|\,dt
= \frac{\|x\|^2}{2\sigma^2},
\]
which is finite.  Thus the degenerate metric produces a well-defined
(pseudo-)distance that assigns the origin finite distance from every
point---geometrically realising the mode as a ``focal point'' of
the distribution.

\paragraph{Connection to score-based models.}
In score-based generative modelling~\cite{SongErmon2019}, the
\emph{score function}
$\nabla \log p$ plays the role of a velocity field.  The gradient
rule converts this velocity field into a Riemannian metric whose
geometry encodes the concentration structure of the distribution.
The degeneration at the mode is the geometric reason that
score-based models require $\varepsilon$-regularisation: the
velocity vanishes at the target, creating a metric singularity.

\begin{corollary}[Finite pseudo-distance to the degeneration locus for gradient metrics]
\label{cor:finite-distance-gradient}
Let $(M,h)$ be connected and let $s\in C^\infty(M)$.
For the gradient rule $g := |\nabla s|_h^2\,h$, the degeneration locus is
$\mathcal{D}(\Phi,s)=\{\nabla s=0\}=\Crit(s)$.  If $\Crit(s)\neq\emptyset$ then
every point of $M\setminus\Crit(s)$ has finite $g$-length to $\Crit(s)$, hence
finite (pseudo-)distance to the degeneration locus.
\end{corollary}

\begin{proof}
Choose $p\in\Crit(s)$ and any piecewise smooth path $\gamma$ from $x$ to $p$.
Since $g=f\,h$ with $f=|\nabla s|_h^2$, we have
$|\gamma'|_g=\sqrt{f(\gamma)}\,|\gamma'|_h=|\nabla s(\gamma)|_h\,|\gamma'|_h$.
The image of $\gamma$ is compact, so $|\nabla s|_h$ is bounded along it and
$L_g(\gamma)<\infty$.
\end{proof}

\begin{corollary}[Stability of the degeneration locus for gradient metrics]
\label{cor:crit-stability}
Let $(M,h)$ be compact without boundary and let $s_0\in C^\infty(M)$ be a
Morse function.  Then there exists a $C^2$-neighbourhood $\mathcal{U}$ of
$s_0$ such that for every $s\in\mathcal{U}$:
\begin{enumerate}[label=\textup{(\alph*)}]
\item $s$ is Morse and $\Crit(s)$ consists of finitely many points.
\item There is a canonical bijection between $\Crit(s_0)$ and $\Crit(s)$
  (each critical point persists uniquely under the perturbation), and
  the Morse index is preserved.
\item For the gradient rule $g_s:=|\nabla s|_h^2\,h$, the degeneration locus
  $\mathcal{D}(\Phi,s)=\Crit(s)$ varies continuously with $s$ (in the
  Hausdorff topology on closed subsets of $M$).
\end{enumerate}
\end{corollary}

\begin{proof}
Since $s_0$ is Morse, each critical point $p\in\Crit(s_0)$ is nondegenerate,
so by the implicit function theorem the equation $\nabla s=0$ has a unique
solution near $p$ for all $s$ sufficiently close to $s_0$ in $C^2$.
Moreover, the Hessian depends continuously on $s$ in $C^2$, so the Morse index
is locally constant and hence preserved.

To rule out creation of new critical points away from $\Crit(s_0)$, pick
disjoint neighbourhoods $U_p$ of each $p\in\Crit(s_0)$ and set
$K:=M\setminus\bigcup_p U_p$.  Then $|\nabla s_0|_h$ attains a positive minimum
on $K$.  For $s$ sufficiently close to $s_0$ in $C^1$, the gradient
$\nabla s$ is uniformly close to $\nabla s_0$ on $K$, so $\nabla s$ cannot
vanish on $K$.  This yields the bijection and the Hausdorff continuity.
\end{proof}

% ----------------------------------------------------------------
\subsection{Warped collapse to a point}
\label{sec:ex-warped}
% ----------------------------------------------------------------

Let $(F, g_F)$ be a compact Riemannian manifold (e.g., $S^{n-1}$
with the round metric), and let $M = [0,1] \times F$ with product
metric $h = dt^2 + g_F$.  Define the seam $s(t,y) = t$ (the
``height function'') and the warped rule with $\phi(s) = s$:
\[
(\alpha, \beta, \gamma) = (s^2, 0, 0), \qquad
g = t^2\, g_F.
\]
(The $dt^2$-component is suppressed: the warped metric acts on the
fibre directions only.)

\paragraph{Geometry near $t=0$.}
At $t = t_0 > 0$, the fibre $(F, t_0^2\, g_F)$ is a rescaled copy
of $F$ with diameter $\diam(F,t_0^2 g_F) = t_0 \cdot \diam(F, g_F)$.
As $t_0 \to 0$, this diameter shrinks to zero: the fibre collapses
to a point.

\paragraph{Gromov--Hausdorff collapse.}
The sublevel set $\{t \le \varepsilon\} \times F$, equipped with
the metric $g_\varepsilon = t^2 g_F + dt^2$, converges in the
Gromov--Hausdorff sense to the interval $[0,\varepsilon]$ as the
fibre shrinks.  This is a concrete instance of the following general
principle: \emph{the degeneration locus of the warped rule
($\phi(s) = 0$ at $s = 0$) is the geometric collapse locus}.

\paragraph{Curvature blowup.}
For the warped product $g = dt^2 + t^2 g_F$ (a metric cone), the
sectional curvature of a plane containing $\partial_t$ is
\[
K(\partial_t, X) = 0,
\]
while a plane tangent to the fibre has sectional curvature
\[
K(X, Y) = \frac{1 - t^2 K_F(X,Y)}{t^2},
\]
where $K_F$ is the sectional curvature of $(F, g_F)$.  As $t \to 0$,
$K(X,Y) \to +\infty$ (unless $K_F \equiv 1$, in which case the cone
is flat).  The curvature blowup at the degeneration locus signals
the formation of a conical singularity.

% ----------------------------------------------------------------
\subsection{A mixed conformal--Hessian geometry}
\label{sec:ex-mixed}
% ----------------------------------------------------------------

To illustrate the power of using multiple terms simultaneously,
consider $M = \mathbb{R}^2$ with $h = \delta$ and the seam
$s(x,y) = x^2 + y^2$.  Define a mixed rule:
\[
(\alpha, \beta, \gamma) = (1, 0, \tfrac{1}{2}), \qquad
g = h + \tfrac{1}{2}\nabla^2 s
  = \begin{pmatrix} 1 & 0 \\ 0 & 1 \end{pmatrix}
    + \begin{pmatrix} 1 & 0 \\ 0 & 1 \end{pmatrix}
  = 2\,\delta.
\]
This is just a rescaled flat metric---unsurprising for a quadratic
seam.  But now take a more interesting seam:
\[
s(x,y) = x^2 + \tfrac{1}{4}y^4.
\]
Then
\[
\nabla^2 s = \begin{pmatrix} 2 & 0 \\ 0 & 3y^2 \end{pmatrix},
\qquad
g = \delta + \tfrac{1}{2}\nabla^2 s
  = \begin{pmatrix} 2 & 0 \\ 0 & 1 + \tfrac{3}{2}y^2 \end{pmatrix}.
\]
The background $\delta$-term ensures positive-definiteness everywhere
(even where $\nabla^2 s$ alone is degenerate at $y = 0$), while the
Hessian term introduces curvature.  The Gauss curvature is
\[
K = -\frac{1}{2\sqrt{\det g}}\left[
  \partial_y\!\left(
    \frac{\partial_y g_{11}}{2\sqrt{\det g}}
  \right)
  + \partial_x\!\left(
    \frac{\partial_x g_{22}}{2\sqrt{\det g}}
  \right)
\right]
= -\frac{3y}{2(2 + 3y^2)^2},
\]
which changes sign across $y = 0$.  The mixed rule has produced a
surface with a line of zero curvature---a saddle transition---from a
single scalar field.

\subsection{Information geometry: Fisher--Rao metric}
\label{sec:ex-fisher}
% ----------------------------------------------------------------

Let $\{p(x|\theta) : \theta \in \Theta \subset \mathbb{R}^n\}$
be an exponential family
$p(x|\theta) = \exp\bigl(\theta \cdot T(x) - A(\theta)\bigr)\mu(dx)$
with log-partition function $A(\theta) = \log\int e^{\theta \cdot T(x)}\mu(dx)$.
The Fisher information metric~\cite{Amari2016} on the parameter
space~$\Theta$ is
\[
g^{\mathrm{Fisher}}_{ij}(\theta)
= \mathbb{E}_{\theta}\!\bigl[\partial_i \log p\;\partial_j \log p\bigr]
= \partial_i \partial_j A(\theta).
\]
This is \emph{exactly the Hessian rule} ($\alpha = \beta = 0$,
$\gamma = 1$) applied to the seam $s = A(\theta)$ on the flat
background $(\Theta, \delta)$.  Strict convexity of~$A$ (a standard
consequence of the exponential family structure) is precisely the
condition that the Hessian rule produces a positive-definite metric.

From the seam-rule perspective:
\begin{itemize}
\item The \textbf{degeneration locus}
  $\mathcal{D} = \{\theta : \det(\nabla^2 A) = 0\}$
  corresponds to degenerate statistical models where the
  family loses local identifiability.
\item The \textbf{Gauss curvature} (for $n = 2$) is computable
  from~\eqref{eq:gauss-general} with $\lambda_i = \partial_i^2 A$.
  For the normal family $A(\mu,\sigma) = \sigma^2/2 + \mu^2/(2\sigma^2)$
  this recovers the classical constant negative curvature of the
  Fisher--Rao metric on the half-plane model.
\item The \textbf{linearisation} $D\Phi_\theta(\dot\theta) = \nabla^2\dot\theta$
  is the third derivative tensor of~$A$, whose kernel (affine functions)
  captures the trivial reparametrisations of the family.
\end{itemize}
Information geometry thus provides a natural, high-profile instance
of the Hessian rule in which every element of the seam-rule
framework---classification, degeneration, linearisation---has a
concrete statistical interpretation.

% ----------------------------------------------------------------
\subsection{The inverse problem: finding the seam}
\label{sec:ex-inverse}
% ----------------------------------------------------------------

The examples above construct geometries \emph{from} a given seam.
In practice, one often faces the reverse: given a target Riemannian
metric~$g$, find a seam~$s$ (and a rule) such that
$g = \Phi(s;\,h)$.  The Universal Decomposition reduces this to a
structured decision tree.

\begin{remark}[Recipe for finding the seam]
\label{rem:seam-recipe}
Let $(M,h)$ be a fixed background and $g$ a target metric.
\begin{enumerate}[label=\textbf{Step \arabic*.}]
\item \textbf{Try the conformal rule.}
  Write $g = e^{2\sigma}\,h$ pointwise.  If $g$ is conformal to~$h$
  (i.e.\ $g_{ij}/h_{ij}$ is a positive scalar function), then
  \[
    s = \tfrac{1}{2}\ln(g_{ij}/h_{ij})
  \]
  is the seam, with $(\alpha,\beta,\gamma) = (e^{2s},0,0)$.
  This always succeeds on surfaces (by the uniformisation theorem,
  every metric is conformal to a constant-curvature background up to
  diffeomorphism) and succeeds in higher dimensions within a fixed
  conformal class.

\item \textbf{Try the Hessian rule.}
  If $g$ is not conformal to~$h$, test whether $g$ is an exact Hessian.
  On a flat background (e.g.\ $h=\delta$ in a coordinate chart), this is
  equivalent to the classical Saint-Venant conditions
  $\partial_i \partial_j g_{kl} = \partial_k \partial_l g_{ij}$.
  On a curved background, commuting covariant derivatives introduces
  the curvature $\Rm(h)$, so the corresponding \emph{generalized}
  Saint-Venant compatibility identities must be used instead.
  If the compatibility conditions hold, solve $\nabla^2 s = g$
  (a Poisson-type system) to obtain~$s$, with
  $(\alpha,\beta,\gamma) = (0,0,1)$.  The solution is unique up to an
  affine function (the kernel of~$\nabla^2$).

\item \textbf{Use the full three-term family.}
  If neither pure rule suffices, write
  $g = \alpha\,h + \beta\,ds\otimes ds + \gamma\,\nabla^2 s$
  and solve the resulting nonlinear PDE for~$s$ (with
  $\alpha,\beta,\gamma$ as coefficient functions).  The linearised
  version of this problem is the normal equation
  $(D\Phi_{s_0})^*(D\Phi_{s_0})\,\delta s
  = (D\Phi_{s_0})^*(g - \Phi(s_0;h))$,
  which is fourth-order elliptic when $\gamma \neq 0$
  (Proposition~\ref{prop:fredholm}) and amenable to iterative
  solvers.  Local surjectivity
  (Proposition~\ref{prop:local-surj}) guarantees that, near any
  non-degenerate background seam, a solution exists.
\end{enumerate}
\end{remark}

We illustrate with a well-known geometry not among the earlier
examples.

\begin{example}[Poincar\'e disk]
\label{ex:poincare}
The Poincar\'e disk $(\mathbb{D}^2, g_P)$ with background
$h = \delta$ (flat Euclidean) has
\[
g_P = \frac{4}{(1-|z|^2)^2}\;\delta.
\]
This is manifestly conformal to~$h$, so \textbf{Step~1} succeeds:
\[
s(z) = \tfrac{1}{2}\ln\frac{4}{(1-|z|^2)^2}
     = \ln\frac{2}{1-|z|^2}.
\]
One verifies:
\begin{itemize}
\item $s(0) = \ln 2$; \; $s \to +\infty$ as $|z|\to 1$
  (the ideal boundary).
\item $\nabla s = 2z/(1-|z|^2)$, which is nowhere zero on
  $\mathbb{D}$: the degeneration locus is empty.
\item The Gaussian curvature $K = -e^{-2s}\Delta s = -1$,
  confirming constant negative curvature.
\end{itemize}
The Hessian rule fails here: $g_P$ is not an exact Hessian
(in Euclidean coordinates, the classical Saint-Venant conditions
$\partial_i\partial_j g_{kl}=\partial_k\partial_l g_{ij}$ are violated
because the conformal factor depends on both $x$ and~$y$ in a
non-separable way).
The conformal rule is the unique pure basis rule that works.
\end{example}

\begin{remark}[What the seam representation buys]
\label{rem:seam-payoff}
The Poincar\'e metric is classical; the value of its seam
$s = \ln(2/(1-|z|^2))$ lies in what the framework makes
\emph{easy} once the seam is known:
\begin{enumerate}[label=(\roman*)]
\item \textbf{Scalar deformations.}
  Perturbing $s \mapsto s + \varepsilon\,\varphi$ for any
  $\varphi \in C^\infty(\mathbb{D})$ gives a valid Riemannian
  metric $g_\varepsilon = e^{2(s+\varepsilon\varphi)}\delta$ for
  every~$\varepsilon$, with no compatibility conditions to check.
  The first-order curvature response is
  $\dot K = -e^{-2s}\Delta\varphi + 2\varphi\,K$---a linear PDE
  in one unknown.
\item \textbf{Curvature as a scalar equation.}
  The Gaussian curvature is $K = -e^{-2s}\Delta s$ (no Christoffel
  symbols, no Riemann tensor).  Prescribing a target curvature
  function~$K_{\mathrm{target}}$ becomes the semilinear PDE
  $\Delta s + K_{\mathrm{target}}\,e^{2s} = 0$.
\item \textbf{Ricci flow as a scalar evolution.}
  The Ricci-flow equation $\partial_t g = -2Kg$ reduces to
  $\partial_t s = e^{-2s}\Delta s$---a single parabolic PDE
  (cf.\ \S\ref{sec:ricci-flow}).
\item \textbf{Automatic discretisation.}
  Given a mesh on~$\mathbb{D}$ with vertices~$\{v_i\}$, the
  discrete conformal rule
  $\ell_s(u,v) = \ell_0(u,v)\,(e^{s(u)}+e^{s(v)})/2$
  produces a weighted graph with one scalar per vertex
  (\S\ref{sec:discrete}).  This has $|V|$ degrees of freedom,
  versus $3|E|$ for a general discrete metric---an order-of-magnitude
  reduction.
\item \textbf{Cross-domain transfer.}
  Every theorem stated for a general seam---the Morse bound on the
  degeneration locus (\S\ref{sec:morse-bound}), the Fredholm
  structure of the inverse problem (\S\ref{sec:ellipticity}),
  the Gauss--Bonnet constraints on the coefficients
  (\S\ref{sec:curvature})---applies verbatim to~$s$.  The seam is
  a common interface: a result proved for one seam
  (e.g.\ the diffusion-model seam
  $\log p$) transfers to the
  Poincar\'e seam and vice versa.
\end{enumerate}
In short, the Poincar\'e metric is old news; the Poincar\'e
\emph{seam} is a computational and theoretical handle that lets one
deform, discretise, flow, and optimise the geometry through a single
scalar field.
\end{remark}

\paragraph{Summary of the examples.}
The seven computations above demonstrate that the Universal
Decomposition is not merely a formal classification: different
choices of $(\alpha, \beta, \gamma)$ produce qualitatively distinct
geometries (constant curvature, Lorentzian signature, metric
degeneration, conical collapse, curvature-sign transitions, Fisher
information, hyperbolic space), all from elementary scalar fields.
The inverse recipe (Remark~\ref{rem:seam-recipe}) shows that
finding the seam for a given metric is a structured problem with
a clear decision tree.

% ================================================================
\section{Curvature of the General Three-Term Metric}
\label{sec:curvature}
% ================================================================

We now derive the scalar curvature of the full three-term metric
$g = \alpha\,h + \beta\,ds \otimes ds + \gamma\,\nabla^2 s$
in terms of the seam and its jets.  This makes precise the sense in
which seam data controls curvature.

Restrict for clarity to $n = 2$ and a flat background $h = \delta$.
Write $\lambda_1 = \alpha + \beta\,(\partial_1 s)^2 + \gamma\,s_{11}$
and $\lambda_2 = \alpha + \beta\,(\partial_2 s)^2 + \gamma\,s_{22}$
for the diagonal components (working in coordinates adapted so that
$s_{12} = 0$ at the point of interest, which can always be arranged
by an $\OO(2)$-rotation of the frame).  Then $\det g =
\lambda_1 \lambda_2$ and the Gauss curvature is
\begin{equation}
\label{eq:gauss-general}
K = -\frac{1}{2\sqrt{\det g}}\left[
  \partial_2\!\left(\frac{\partial_2 \lambda_1}{\sqrt{\det g}}\right)
  + \partial_1\!\left(\frac{\partial_1 \lambda_2}{\sqrt{\det g}}\right)
\right].
\end{equation}
Expanding, one finds that $K$ is a rational function of the
$\OO(n)$-invariants $\mathcal{I}$ (equation~\eqref{eq:invariants})
and the derivatives of $\alpha, \beta, \gamma$.

\begin{remark}[Structure of the curvature formula]
The key structural observation is:
\begin{itemize}
\item \textbf{Zeroth-order contributions} come from the $\alpha$-term
  and produce conformal curvature $K_{\text{conf}} = -(\Delta\log\alpha)/(2\alpha)$.
\item \textbf{First-order contributions} come from the
  $\beta$-term and couple gradient invariants to curvature.
\item \textbf{Second-order contributions} come from the
  $\gamma$-term and involve derivatives of $s$ of order at most~$3$.
  Although differentiating the Christoffel symbols of
  $g=\alpha h+\beta\,ds\otimes ds+\gamma\,\nabla^2 s$ naively produces
  fourth derivatives of $s$, these cancel in the alternating
  combination defining the Riemann tensor (a Hessian-geometry
  cancellation).
\end{itemize}
Thus the scalar curvature of a general three-term metric is, in
general, a \emph{third-order} nonlinear expression in the seam.
\end{remark}

\subsection{General dimension and curved background}
\label{sec:curv-general}

The explicit formula~\eqref{eq:gauss-general} was stated for $n=2$
and flat~$h$.  We now describe the structure for arbitrary~$n$ and
a curved background $(M^n,h)$.

\begin{proposition}[General curvature structure]
\label{prop:curv-general}
Let $g = \alpha\,h + \beta\,ds\otimes ds + \gamma\,\nabla^2 s$
be a Riemannian metric on $(M^n,h)$ with $\alpha,\beta,\gamma$
depending on $(s, |\nabla s|^2)$.  Then the Riemann curvature of
$g$ (viewed as the $(1,3)$ curvature tensor) has the form
\begin{equation}
\label{eq:Riem-structure}
\Rm(g) = \Rm(h)
+ \mathcal{C}_1(\alpha,\beta,\gamma,s;\nabla s,\nabla^2 s)
+ \mathcal{C}_2(\alpha,\beta,\gamma,s;\nabla^3 s),
\end{equation}
where:
\begin{enumerate}[label=(\roman*)]
\item $\Rm(h)$ is the curvature inherited from the background.
  (Lowering indices and scalar contractions are performed with $g$;
  for instance, for $g=e^{2\varphi}h$ this yields the usual conformal
  scaling in the scalar curvature.)
\item $\mathcal{C}_1$ is a curvature-type tensor that is polynomial
  in $\nabla s$ and $\nabla^2 s$ and rational in
  $\alpha, \beta, \gamma$; it captures the contributions from
  the gradient and Hessian terms.
\item $\mathcal{C}_2$ is nonzero only when $\gamma \neq 0$
  and involves derivatives of $s$ of order $3$; it arises
  because the Christoffel symbols of $g$ contain $\nabla^3 s$
  via the $\gamma\,\nabla^2 s$ term, and the would-be $\nabla^4 s$
  contributions cancel in the alternating combination defining $\Rm(g)$.
\end{enumerate}
In particular, the scalar curvature $R(g) = g^{ij}g^{kl}R_{ikjl}$
is a rational function of $\alpha,\beta,\gamma$, the
$\OO(n)$-invariants~$\mathcal{I}$
(equation~\eqref{eq:invariants}), and the background curvature
$\Rm(h)$, involving derivatives of $s$ up to order~$3$.
\end{proposition}

\begin{proof}[Proof sketch]
Write $g = \alpha\,h + P$ where $P = \beta\,ds\otimes ds +
\gamma\,\nabla^2 s$.  The Christoffel symbols of $g$ relative to
those of $h$ are given by the Koszul formula:
$\Gamma(g) - \Gamma(h) = \tfrac{1}{2}g^{-1}(\nabla^h P +
\nabla^h P - \nabla^h P)$ (with appropriate index symmetrisations).
Since $P$ involves $\nabla s$ and $\nabla^2 s$, the connection
difference involves $\nabla^2 s$ and $\nabla^3 s$.
Differentiating once more for the Riemann tensor and contracting
twice for the scalar curvature gives the stated structure.
The key point is that $g^{-1}$ introduces $\det(g)$ in the
denominator, making $R(g)$ rational in the coefficients.
\end{proof}

\begin{remark}[Reduction to known formulas]
For $\beta = \gamma = 0$, $\alpha = e^{2\varphi}$:
$R(g) = e^{-2\varphi}\bigl[R(h) - 2(n-1)\Delta_h\varphi
- (n-1)(n-2)|\nabla\varphi|^2_h\bigr]$,
recovering the standard conformal change formula.
For $\alpha = \beta = 0$, $\gamma = 1$ on $(M^2,\delta)$:
$R(g) = 2K$ with $K$ given by~\eqref{eq:gauss-general},
recovering the two-dimensional Hessian result.
\end{remark}

\subsection{Gauss--Bonnet constraints on the coefficients}
\label{sec:gauss-bonnet}

On a closed oriented surface $(M^2, h)$, the Gauss--Bonnet theorem
provides a global constraint on any seam-generated metric.

\begin{proposition}[Gauss--Bonnet for seam metrics]
\label{prop:gauss-bonnet}
Let $g = \alpha\,h + \beta\,ds\otimes ds + \gamma\,\nabla^2 s$
be a Riemannian metric on a closed surface $(M^2, h)$.
Then
\[
\int_M K_g\,dA_g = 2\pi\chi(M),
\]
where $K_g$ is the Gauss curvature of~$g$
and $dA_g = \sqrt{\det g}\,dx^1\wedge dx^2$.  In particular,
the seam $s$ and coefficients $\alpha,\beta,\gamma$ must satisfy
the integral constraint
\[
\int_M K_g(\alpha,\beta,\gamma,s,\nabla s,\nabla^2 s,\nabla^3 s)
\;\sqrt{\det g(\alpha,\beta,\gamma,s,\nabla s,\nabla^2 s)}
\;dA_h = 2\pi\chi(M).
\]
\end{proposition}

\begin{corollary}[Weighted Gauss--Bonnet for order-1 rules]
\label{cor:weighted-gb-order1}
Let $g=\alpha\,h+\beta\,ds\otimes ds$ be an order-$1$ seam metric on a closed
surface $(M^2,h)$, with $\alpha>0$ and $\alpha+\beta|\nabla s|_h^2>0$.
Then
\[
\int_M K_g\;\alpha\,\sqrt{1+\frac{\beta}{\alpha}|\nabla s|_h^2}\;dA_h
= 2\pi\chi(M).
\]
\end{corollary}

\begin{proof}
By Gauss--Bonnet, $\int_M K_g\,dA_g=2\pi\chi(M)$.  Since $\dim M=2$,
Corollary~\ref{cor:volume-order1} gives
$dA_g=dV_g=\alpha\,\sqrt{1+\frac{\beta}{\alpha}|\nabla s|_h^2}\,dA_h$.
Substituting this expression for $dA_g$ yields the identity.
\end{proof}

This constrains the admissible parameter space in concrete ways:
for instance, a seam on the 2-sphere ($\chi = 2$) generating a metric
via the conformal rule must produce strictly positive total curvature,
while on a torus ($\chi = 0$) the positive and negative curvature
contributions from the $\beta$- and $\gamma$-terms must cancel exactly.

\begin{remark}[Chern--Gauss--Bonnet in dimension~4]
\label{rem:4d-cgb}
In dimension~4, the Chern--Gauss--Bonnet integrand
$|\Rm|^2 - 4|\Ric|^2 + R^2$ can be expressed, via
Proposition~\ref{prop:curv-general}, entirely in terms of jets
of~$s$ up to order~3 and the coefficient functions
$(\alpha,\beta,\gamma)$.  Integrating against $dV_g$ yields a
\emph{seam functional}
\[
\chi(M) = \frac{1}{8\pi^2}
\int_M Q\bigl(s,\nabla s,\nabla^2 s,\nabla^3 s;\,
\alpha,\beta,\gamma\bigr)\,dV_h,
\]
where $Q$ is an explicit polynomial in the jet invariants.
Since $\chi(M)$ is a fixed integer, this constrains the admissible
seams on any given 4-manifold: a nonlinear integral identity that
every metric-generating seam must satisfy.  For the pure Hessian
rule ($\alpha=\beta=0$, $\gamma=1$) on a flat background,
$Q$ simplifies to a polynomial in the eigenvalues of $\nabla^2 s$
and their derivatives, giving a new route to seam selection
(Open Problem~\ref{prob:seam-selection}).
\end{remark}

\subsection{Local scalar-curvature prescription}
\label{sec:curv-prescription}

The curvature formula (Proposition~\ref{prop:curv-general}) combined
with the Fredholm theory for the normal equation
(Proposition~\ref{prop:fredholm}) suggests an avenue for inverse design.
However, the scalar-curvature map itself does \emph{not} inherit a
fourth-order elliptic principal symbol in the seam variable: the would-be
$|\xi|^4$ term cancels.

\begin{proposition}[Degeneracy of scalar-curvature prescription]
\label{prop:curv-prescription}
Let $(M^n,h)$ be compact without boundary and let $s_0$ be a smooth seam.
Consider the scalar-curvature map
$\mathcal{R}\colon s\mapsto R(\Phi(s;h))$.
Then the linearisation $D\mathcal{R}_{s_0}(\dot s)$ has \emph{no}
fourth-order term in $\dot s$ coming from the Hessian direction
$\gamma\,\nabla^2\dot s$: the $|\xi|^4$ principal symbol cancels.
In particular, $D\mathcal{R}_{s_0}$ is not, in general, a fourth-order
elliptic scalar operator in the seam variable.

On a closed surface ($n=2$), any prescribed target curvature must still
satisfy the Gauss--Bonnet constraint
$\int_M R_{\mathrm{target}}\,dA_g = 4\pi\chi(M)$.
\end{proposition}

\begin{proof}
The linearisation of scalar curvature at $g_0=\Phi(s_0;h)$ in the
direction of a symmetric 2-tensor $\dot g$ is
\[
\delta R = -\Delta_{g_0}(\tr_{g_0} \dot g) + \nabla^{i}\nabla^{j} \dot g_{ij}
- \langle \Ric(g_0), \dot g\rangle_{g_0}.
\]
If we take the Hessian-direction variation $\dot g=\gamma\,\nabla^2\dot s$
then the symbol-level contribution of the first two terms cancels:
in a normal frame at a point,
$\sigma(-\Delta\,\tr)(\xi)=-\gamma|\xi|^4$ while
$\sigma(\nabla^i\nabla^j \dot g_{ij})(\xi)=+\gamma|\xi|^4$.
Hence the $|\xi|^4$ principal symbol vanishes.
The remaining terms involve at most lower derivatives of $\dot s$.
\end{proof}

\begin{remark}[Connection to Kazdan--Warner]
On closed surfaces with the conformal rule
($\beta = \gamma = 0$, $\alpha = e^{2s}$),
the scalar-curvature prescription problem reduces to the
classical Kazdan--Warner theorem~\cite{KazdanWarner1975}: a function
$K_{\mathrm{target}}$ that is realised as the Gaussian curvature of a
conformal metric $e^{2s}h$ must satisfy the Gauss--Bonnet constraint
$\int_M K_{\mathrm{target}}\,e^{2s}\,dA_h = 2\pi\chi(M)$.
The Kazdan--Warner theorem characterises which curvature functions are
realisable beyond this necessary integral condition.
For the full three-term family, Proposition~\ref{prop:curv-prescription}
shows that the seam-linearised scalar curvature does \emph{not} have a
fourth-order elliptic principal symbol coming from the Hessian direction,
so a general local prescription theorem requires additional structure.
\end{remark}

\subsection{The seam Ricci flow}
\label{sec:ricci-flow}

A time-dependent seam $s(x,t)$ induces an evolving metric
$g(t) = \Phi(s(\cdot,t);\,h)$.  Requiring that $g(t)$ satisfy the
Ricci flow $\partial_t g = -2\Ric(g)$ and using the chain rule
$\partial_t g = D\Phi_s(\partial_t s)$ yields the \emph{seam Ricci
flow equation}:
\begin{equation}
\label{eq:seam-ricci}
D\Phi_s(\partial_t s) = -2\,\Ric(\Phi(s;\,h)).
\end{equation}
For the conformal rule ($\beta = \gamma = 0$,
$\alpha = e^{2s}$, $n = 2$), the linearisation is
$D\Phi_s(\dot s) = 2\dot s\,g$ and the Ricci tensor is
$\Ric(g) = K\,g$, so~\eqref{eq:seam-ricci} reduces to
$\partial_t s = -K = e^{-2s}\Delta s$---the standard equation
for the Ricci flow on surfaces in conformal coordinates.

For the Hessian rule ($\alpha = \beta = 0$, $\gamma = 1$),
the curvature of $g = \nabla^2 s$ is controlled by the third jet
$\nabla^3 s$ (cf.\ \S\ref{sec:curv-general}).  A convenient strictly
elliptic scalar reduction of the seam Ricci flow is obtained by applying
the normal-operator reduction in
Proposition~\ref{prop:ricci-wellposed} below, yielding a \emph{fourth-order}
scalar evolution equation in which the term containing $\partial_t s$ is
governed by a fourth-order elliptic operator.
This is the seam-geometric avatar of the Ricci flow, reducing
the evolution of an entire metric to the evolution of a single
scalar field.

\begin{remark}[Connection to Calabi flow]
\label{rem:calabi}
On a K\"ahler manifold $(M, \omega, J)$ with K\"ahler potential
$\varphi$ (so that $\omega = dd^c \varphi$), the \emph{Calabi flow}
is $\partial_t \varphi = R(\omega_{\varphi})$---a fourth-order
parabolic PDE on a scalar potential whose solution evolves
the entire K\"ahler metric.  The Hessian-rule seam Ricci flow
(with $\gamma = 1$, $\alpha = \beta = 0$) is a real analogue in the
sense that one can reduce the metric evolution to a fourth-order scalar
evolution equation for a potential; in both cases, a fourth-order scalar
equation arises after an appropriate reduction.
The key difference is that Calabi flow preserves the K\"ahler
class, while the seam Ricci flow operates without a complex
structure and preserves only the seam representation itself.
\end{remark}

\begin{proposition}[Well-posedness of the seam Ricci flow]
\label{prop:ricci-wellposed}
Let $(M^n,h)$ be compact without boundary and
$\Phi = \alpha\,h + \beta\,ds\otimes ds + \gamma\,\nabla^2 s$
with smooth coefficients depending on $(s,|\nabla s|^2)$.
\begin{enumerate}[label=(\roman*)]
\item \textbf{Strictly parabolic regime ($\gamma > 0$):}
  Fix an initial seam $s_0 \in C^\infty(M)$ such that
  $g_0 = \Phi(s_0;\,h)$ is a Riemannian metric.
  Then the normal equation
  $(D\Phi_{s_0})^*\bigl[D\Phi_s(\partial_t s) + 2\Ric(g)\bigr] = 0$
  is a quasilinear evolution equation in which the term containing
  $\partial_t s$ is governed by a fourth-order uniformly elliptic
  operator.  More precisely, after freezing coefficients, its
  principal part in $\partial_t s$ is $\gamma(s_0)^2\,\Delta_h^2\,\partial_t s$
  (where $\Delta_h^2$ denotes the bi-Laplacian with respect to~$h$).
  Consequently, after fixing a standard normalisation to remove the
  finite-dimensional kernel of the normal operator (e.g.\ prescribing
  the mean value of $s$), there exists $T > 0$ and a unique smooth
  solution $s(\cdot,t)$ for $t \in [0,T)$.
\item \textbf{Degenerate regime ($\gamma = 0$):}
  When $\gamma = 0$, the conformal rule
  $\Phi = \alpha\,h + \beta\,ds\otimes ds$ gives a
  \emph{second-order} seam Ricci flow, but the linearised operator
  has zero principal symbol whenever $\nabla s_0 = 0$.
  The flow is well-posed only away from the degeneration locus
  $\{\nabla s = 0\}$, where the parabolicity degenerates.
\end{enumerate}
\end{proposition}

\begin{proof}
\emph{Regime (i): $\gamma>0$.}
Write $g(t)=\Phi(s(\cdot,t);h)$ and consider the normal equation
\begin{equation}
\label{eq:normal-seam-ricci}
(D\Phi_{s_0})^*\bigl[D\Phi_s(\partial_t s) + 2\Ric(g)\bigr] = 0.
\end{equation}
We first identify the principal part of the term containing
$\partial_t s$.

Since the coefficients depend only on $(s,|\nabla s|^2)$, the
linearisation $D\Phi_s$ has leading term $\gamma(s,|\nabla s|^2)\,\nabla_h^2$
acting on the seam variation; all other contributions are of order at
most~$1$ in the seam variation (coming from differentiating
$\alpha h$ and $\beta\,ds\otimes ds$).
Thus, in principal symbol,
\[
\sigma\bigl(D\Phi_s\bigr)(\xi)\cdot u = \gamma(s)\,(\xi\otimes\xi)\,u.
\]
By Proposition~\ref{prop:fredholm}, the formal adjoint
$(D\Phi_{s_0})^*$ is a second-order scalar operator with principal symbol
\[
\sigma\bigl((D\Phi_{s_0})^*\bigr)(\xi)\cdot T
= \gamma(s_0)\,\langle T,\xi\otimes\xi\rangle.
\]
Composing these symbols gives
\[
\sigma\bigl((D\Phi_{s_0})^*\circ D\Phi_s\bigr)(\xi)\cdot u
= \gamma(s_0)\,\gamma(s)\,|\xi|_h^4\,u,
\]
so, in particular at $t=0$, the principal symbol governing
$\partial_t s$ is $\gamma(s_0)^2|\xi|_h^4$.  Equivalently, the frozen
principal part of~\eqref{eq:normal-seam-ricci} is
$\gamma(s_0)^2\,\Delta_h^2\,\partial_t s$.
Since $\gamma(s_0)>0$, this yields a uniformly strongly elliptic operator
in the $\partial_t s$ term.

To obtain short-time existence, rewrite~\eqref{eq:normal-seam-ricci} as
\begin{equation}
\label{eq:normal-seam-ricci-solve}
\mathcal{N}_{s_0,s}(\partial_t s)
= -\,(D\Phi_{s_0})^*\bigl(2\Ric(\Phi(s;h))\bigr),
\end{equation}
where $\mathcal{N}_{s_0,s}:=(D\Phi_{s_0})^*\circ D\Phi_s$ is a fourth-order
scalar operator whose principal symbol is $\gamma(s_0)\gamma(s)|\xi|^4$.
For $s$ sufficiently close to $s_0$ in a high Sobolev norm, ellipticity is
uniform and the standard elliptic estimates imply that the right-hand side
of~\eqref{eq:normal-seam-ricci-solve} determines $\partial_t s$ uniquely
after fixing a normalisation that removes the finite-dimensional kernel of
the normal operator (for instance, by prescribing $\int_M s\,dV_h$).
Since $s\mapsto \Ric(\Phi(s;h))$ depends smoothly on the jet of $s$ and
$\mathcal{N}_{s_0,s}$ depends smoothly on $(s,\nabla s)$, the resulting
right-hand side defines a smooth vector field on a Sobolev manifold of
seams.  Standard local well-posedness for such quasilinear evolution
equations gives a time $T>0$ and a unique smooth solution with initial data
$s(\cdot,0)=s_0$.

\emph{Regime (ii): $\gamma=0$.}
When $\gamma=0$, the highest-order term in $D\Phi_{s_0}(u)$ comes from the
variation of $\beta\,ds\otimes ds$ and is first order in~$u$.
In particular, in principal symbol,
\[
\sigma(D\Phi_{s_0})(\xi)\cdot u
= 2\beta\,(\nabla s_0\cdot\xi)\,(\xi\odot\nabla s_0^\flat)\,u,
\]
which vanishes when $\nabla s_0=0$ and also in directions $\xi\perp\nabla s_0$.
Therefore strict parabolicity/ellipticity fails at critical points and in
transverse directions, so the flow degenerates along $\{\nabla s=0\}$.
\end{proof}

% ================================================================
\section{Beyond Metrics: Quadratic Differentials as a Complex Rule}
\label{sec:qd}
% ================================================================

On a Riemann surface $(M^2, h, J)$ with complex structure $J$, the
target bundle $\mathcal{G}$ may be the bundle of holomorphic quadratic
differentials $K_M^{\otimes 2}$ (holomorphic sections of
$\Sym^2(T^{*1,0}M)$).  Given a seam $s$ and a holomorphic
function $f$ with $\operatorname{Re} f = s$ (i.e.\ $s$ is the real
part of a holomorphic potential), the construction
\begin{equation}
\label{eq:qd-rule}
w\,dz^2 := (\partial_z^2 f)\,dz^2 = f''(z)\,dz^2
\end{equation}
defines a holomorphic quadratic differential away from the zeros of the
underlying function.

\begin{proposition}[QD rule as a complex Hessian]
The quadratic-differential rule~\eqref{eq:qd-rule} is the
$(2,0)$-projection of the complexified Hessian rule.  Specifically,
if $g_{ij} = \partial_i \partial_j s$ is the real Hessian and $J$ is
the complex structure, then
\[
w = \tfrac{1}{2}\bigl(g_{11} - g_{22}\bigr) - i\, g_{12}
\quad\text{(in conformal coordinates $z = x^1 + ix^2$)},
\]
and $g = \operatorname{Re}(w\,dz^2) + \tfrac{1}{2}(\Delta s)\,|dz|^2$.
\end{proposition}

This unifies the Hessian rule (real part) with the
quadratic-differential construction that arises in complex
analysis and mathematical physics.  The
traceless part of the Hessian \emph{is} the quadratic differential;
the trace part is the conformal factor $\Delta s$.

% ================================================================
\section{The Linearisation of a Rule}
\label{sec:linearisation}
% ================================================================

Given a rule $\Phi$ and a background seam $s_0$, the first-order
response of the geometry to a seam perturbation $\dot s$ is
\begin{equation}
\label{eq:linearisation}
D\Phi_{s_0}(\dot s) :=
\left.\frac{d}{dt}\right|_{t=0} \Phi(s_0 + t \dot s;\, h).
\end{equation}
This is a \emph{linear differential operator}
$D\Phi_{s_0} \colon C^\infty(M) \to \Gamma(\Sym^2 T^*M)$ whose
properties---kernel, image, ellipticity, spectrum---control the
geometry of the inverse problem.

\subsection{The linearised operator for each basis rule}

For the universal decomposition
$\Phi = \alpha\,h + \beta\,ds \otimes ds + \gamma\,\nabla^2 s$
with coefficients depending on $(s, |\nabla s|^2, \Delta s)$:

\begin{proposition}[Linearisation at $s_0 = 0$ with $\nabla s_0 = 0$]
\label{prop:lin}
Assume $(M,h)$ is compact without boundary and
$s_0 \equiv 0$ (so $\nabla s_0 = 0$ and $\nabla^2 s_0 = 0$).
Then the linearisation of each basis rule at $s_0 = 0$ is:
\begin{enumerate}[label=(\alph*)]
\item \textbf{Conformal:} $\Phi = e^{2s} h$.
  $D\Phi_0(\dot s) = 2\dot s\, h$.
  This is a zeroth-order (multiplication) operator.
  Kernel: $\dot s = 0$.  Image: conformal deformations.
\item \textbf{Gradient tensor:} $\Phi = |\nabla s|^2 h$.
  $D\Phi_0(\dot s) = 0$.
  The linearisation at $s_0 = 0$ vanishes identically---the gradient
  rule is \emph{degenerate at constant seams}.
\item \textbf{Hessian:} $\Phi = \nabla^2 s$.
  $D\Phi_0(\dot s) = \nabla^2 \dot s$.
  This is a second-order elliptic operator (the Hessian operator).
  Kernel (on compact $M$): affine functions, i.e.\
  $\dot s \in \ker\nabla^2 = \{c_0 + c_i x^i\}$ (dimension $\le n+1$,
  with equality on $\mathbb{R}^n$).  Image:
  all symmetric 2-tensors of the form $\nabla^2 f$ for some~$f$.
\end{enumerate}
\end{proposition}

\begin{proof}
Direct differentiation of each rule.  For (b), the linearisation of
$|\nabla s|^2 h$ at $s_0 = 0$ is
$2\langle \nabla s_0, \nabla \dot s \rangle h = 0$ since
$\nabla s_0 = 0$.  For~(c), $\nabla^2(s_0 + t\dot s)
= \nabla^2 s_0 + t\nabla^2\dot s$, so the linearisation is
$\nabla^2 \dot s$.
\end{proof}

\subsection{The seam deformation complex}

For a given rule $\Phi$ and seam $s_0$, the linearisation
$D\Phi_{s_0}$ fits into a sequence:
\begin{equation}
\label{eq:complex}
0 \longrightarrow
\underbrace{\ker D\Phi_{s_0}}_{\text{gauge / kernel}}
\longrightarrow
C^\infty(M)
\xrightarrow{\;D\Phi_{s_0}\;}
\Gamma(\Sym^2 T^*M)
\longrightarrow
\underbrace{\coker D\Phi_{s_0}}_{\text{obstructions}}
\longrightarrow 0.
\end{equation}

\begin{definition}[Seam moduli space]
The (formal, infinitesimal) \emph{seam moduli space} at $s_0$
for target geometry $g_0 = \Phi(s_0)$ is
$\ker D\Phi_{s_0}$---the space of infinitesimal seam variations
that leave the geometry unchanged to first order.
\end{definition}

\begin{definition}[Obstruction space]
The \emph{obstruction space} at $s_0$ is $\coker D\Phi_{s_0}$---the
space of infinitesimal geometric deformations that cannot be achieved
by any seam variation through the rule~$\Phi$.
\end{definition}

For the conformal rule, the kernel is trivial (0-dimensional) and
the obstruction space is the space of non-conformal deformations
($(n(n+1)/2 - 1)$-dimensional at each point)---consistent with
the well-known fact that a single scalar controls only one
component of the metric.

For the Hessian rule on $\mathbb{R}^n$, the kernel is $(n+1)$-dimensional
(affine functions), and the image consists of exact Hessians
$\nabla^2 f$, which form a proper subspace of all symmetric 2-tensors
(on $\mathbb{R}^n$, a symmetric 2-tensor is an exact Hessian if and only
if it satisfies the classical Saint-Venant compatibility conditions
$\partial_i \partial_j g_{kl}=\partial_k \partial_l g_{ij}$).

\begin{remark}[Connection to the cotangent Laplacian]
In the discrete setting of conformal graph metrics, the linearisation
of the angle-defect curvature map with respect to the seam at $s = 0$
is exactly the cotangent Laplacian (see \S\ref{sec:discrete}).  This is
a discrete avatar of the linearised conformal rule: the cotangent
Laplacian is the discrete analogue of the second-order operator
governing the first-order response of curvature to a conformal
perturbation.
\end{remark}

\subsection{Ellipticity and analytical properties}
\label{sec:ellipticity}

For the general three-term rule
$\Phi = \alpha\,h + \beta\,ds\otimes ds + \gamma\,\nabla^2 s$
at background seam~$s_0$, the linearisation
$D\Phi_{s_0}\colon C^\infty(M) \to \Gamma(\Sym^2 T^*M)$
is a second-order linear differential operator whose analytical
properties are governed by its principal symbol.

\begin{proposition}[Principal symbol of the full linearised operator]
\label{prop:symbol}
At a point $x \in M$ and covector $\xi \in T^*_xM$,
the principal symbol of~$D\Phi_{s_0}$ is the linear map
$\sigma(D\Phi_{s_0})(\xi)\colon \mathbb{R} \to \Sym^2 T^*_xM$
given by
\begin{equation}
\label{eq:symbol}
\sigma(D\Phi_{s_0})(\xi)\cdot\dot s
= \bigl[\gamma\,\xi \otimes \xi
  + 2\beta\,(\nabla s_0 \cdot \xi)\,\xi \otimes \nabla s_0^\flat
  + \text{lower-order}\bigr]\,\dot s,
\end{equation}
where the ``lower-order'' terms come from derivatives of the
coefficient functions $\alpha,\beta,\gamma$ with respect to
$\Delta s$ and $|\nabla s|^2$, and contribute only when
these coefficients genuinely depend on second-order invariants.
\end{proposition}

\begin{proof}
The highest-order part of $D\Phi_{s_0}(\dot s)$ comes from
differentiating $\gamma\,\nabla^2 s$ and from the variation of
$\beta\,ds\otimes ds$ through the dependence of $\beta$ on
$\Delta s$.  The $\nabla^2\dot s$ term contributes
$\gamma\,\xi\otimes\xi$; the cross term from
$\partial_{\Delta s}\beta \cdot \Delta\dot s \cdot ds_0\otimes ds_0$
contributes $(\partial_{\Delta s}\beta)\,|\xi|^2\,\nabla s_0^\flat
\otimes \nabla s_0^\flat$.
For coefficients depending only on $(s, |\nabla s|^2)$---which
covers all worked examples in \S\ref{sec:worked}---the
leading symbol is exactly $\gamma\,\xi\otimes\xi$.
\end{proof}

\begin{corollary}[Ellipticity criterion]
\label{cor:elliptic}
When the coefficients $\alpha,\beta,\gamma$ depend only on
$(s, |\nabla s|^2)$, the linearised operator $D\Phi_{s_0}$ is
\emph{overdetermined-elliptic} (injective symbol) if and only if
$\gamma(s_0) \neq 0$.  The symbol is injective iff
$\sigma(\xi) \neq 0$ for all $\xi \neq 0$: since
$\gamma\,\xi\otimes\xi = 0$ requires $\xi = 0$ when $\gamma \neq 0$,
injectivity holds.
\end{corollary}

The linearised operator maps sections of a rank-1 bundle to sections
of a rank-$n(n+1)/2$ bundle, so it is necessarily overdetermined for
$n \geq 2$.  In the overdetermined-elliptic case, the relevant
Fredholm theory is:

\begin{proposition}[Fredholm property]
\label{prop:fredholm}
Let $(M,h)$ be compact without boundary, $\gamma \neq 0$, and
suppose the coefficients depend only on $(s, |\nabla s|^2)$.
Then the linearisation $D\Phi_{s_0}\colon H^{k+2}(M) \to
H^k(M;\,\Sym^2 T^*M)$ is a bounded linear operator with
\emph{finite-dimensional kernel} and \emph{closed range}.
The formal $L^2$-adjoint $(D\Phi_{s_0})^*\colon
\Gamma(\Sym^2 T^*M) \to C^\infty(M)$ is a second-order operator
with principal symbol
$\sigma((D\Phi_{s_0})^*)(\xi)\cdot T =
\gamma \langle T, \xi\otimes\xi \rangle$
for $T \in \Sym^2 T^*_x M$, and the composition
$(D\Phi_{s_0})^* \circ D\Phi_{s_0}$ is a fourth-order elliptic
scalar operator with principal symbol $\gamma^2 |\xi|^4$.
In particular, the normal operator $(D\Phi_{s_0})^* D\Phi_{s_0}$
is Fredholm of index zero.
\end{proposition}

The vanishing of $\gamma$ is thus the analytical avatar of the
degeneration locus: exactly where the Hessian component of the
rule vanishes, the linearisation loses ellipticity.

% ================================================================
\section{The Degeneration Locus}
\label{sec:degeneration}
% ================================================================

A striking feature revealed by the linearisation is the existence of a
\emph{degeneration locus}: the set of points where the linearised rule
drops rank, meaning that small seam perturbations fail to control the
geometry.

\begin{definition}[Degeneration locus]
\label{def:degen}
For a rule $\Phi$ and a seam $s$, the \emph{degeneration locus} is
\begin{equation}
\mathcal{D}(\Phi, s) :=
\{x \in M : \text{the fibre map }
D\Phi_s(x) \colon J^k_x \to (\Sym^2 T^*M)_x
\text{ is not surjective}\}.
\end{equation}
\end{definition}

This abstract definition specialises to concrete geometric objects
across a range of applications:

\begin{enumerate}[label=(\arabic*)]
\item \textbf{Gradient rule, diffusion seam:}
  $\Phi = |\nabla s|^2 h$.  The degeneration locus is
  $\{\nabla s = 0\}$---the critical points of the log-density, i.e.\
  the modes and saddles of the probability distribution.  This is
  precisely where the probability-flow ODE velocity vanishes and the
  seam metric becomes degenerate, requiring $\varepsilon$-regularisation.

\item \textbf{Hessian rule, convex potential:}
  $\Phi = \nabla^2 s$ for a convex function~$s$.  The degeneration
  locus is $\{\det(\nabla^2 s) = 0\}$, which coincides with the
  singular support of the optimal-transport map in the Brenier--McCann
  theory.  Caffarelli's regularity theory governs the geometry here.

\item \textbf{Warped rule:}
  $\Phi = \phi(s)^2 h$ with $\phi(0) = 0$.  The degeneration
  locus $\{s = 0\}$ is the collapse locus: the fibre shrinks to a
  point and the Gromov--Hausdorff limit along a degeneration
  sequence loses the fibre directions entirely
  (cf.\ Example~\ref{sec:ex-warped}).
\end{enumerate}

\begin{proposition}[The degeneration locus is a geometric invariant]
\label{prop:degen-invariant}
If $\Phi$ is a natural operator (satisfying the naturality axiom
\ref{item:naturality}), then for any diffeomorphism $\varphi$ of $M$
preserving $\mathcal{B}$,
\[
\mathcal{D}(\Phi, \varphi^* s) = \varphi^{-1}(\mathcal{D}(\Phi, s)).
\]
In particular, the degeneration locus is an intrinsic geometric object,
not an artefact of coordinates.
\end{proposition}

\begin{proof}
By naturality, $D\Phi_{\varphi^* s}(\cdot)
= \varphi^* \circ D\Phi_s(\cdot) \circ (\varphi^{-1})^*$.
Surjectivity is preserved by pullback, so a point $x$ lies in
$\mathcal{D}(\Phi, \varphi^* s)$ if and only if $\varphi(x)$ lies in
$\mathcal{D}(\Phi, s)$.
\end{proof}

\begin{remark}[A unifying theme]
Across the examples above, \emph{the most interesting geometry
happens at or near the degeneration locus}: modes of the score
field generate the critical-point regularisation problem;
singular optimal-transport supports delimit the regularity
boundary; collapse loci of warped products mark the loss of
fibre directions.  The degeneration locus is a universal
geometric concept that connects these phenomena.
\end{remark}

\subsection{Morse-theoretic lower bound}
\label{sec:morse-bound}

For rules that depend on the gradient of the seam, the topology of the
manifold forces the degeneration locus to be nonempty.

\begin{theorem}[Topological lower bound on the degeneration locus]
\label{thm:morse-degen}
Let $(M^n, h)$ be a closed Riemannian manifold and let $\Phi$ be a
rule whose linearised symbol $\sigma(D\Phi_s)(\xi)$ vanishes whenever
$\nabla s = 0$\textup{---}in particular, any rule with $\gamma = 0$
and $\alpha$ depending on $|\nabla s|^2$, or more generally any rule
dominated by the gradient tensor.  Then for any Morse seam
$s\colon M \to \mathbb{R}$,
\[
\#\,\mathcal{D}(\Phi, s) \;\geq\; \sum_{k=0}^{n} b_k(M),
\]
where $b_k(M)$ are the Betti numbers of~$M$.  If $s$ is merely
smooth (not Morse), $\mathcal{D}(\Phi,s) \supseteq \operatorname{Crit}(s)$
and the Ljusternik--Schnirelmann category gives
$\#\,\mathcal{D}(\Phi,s) \geq \operatorname{cat}(M) \geq 1 + \operatorname{cup-length}(M)$.
\end{theorem}

\begin{proof}
By hypothesis, $\nabla s(x) = 0$ implies $\sigma(D\Phi_s)(\xi)|_x = 0$
for all~$\xi$, so $x \in \mathcal{D}(\Phi,s)$.  Thus
$\operatorname{Crit}(s) \subseteq \mathcal{D}(\Phi,s)$.
If $s$ is Morse, each critical point is isolated, and the Morse
inequalities give
$\#\,\operatorname{Crit}(s) \geq \sum b_k(M)$.
The Ljusternik--Schnirelmann bound requires only smoothness.
\end{proof}

\begin{corollary}[Poincar\'e--Hopf charge of isolated critical points]
\label{cor:poincare-hopf-degen}
In the setting of Theorem~\ref{thm:morse-degen}, suppose that $\Crit(s)$ is
discrete (for instance, $s$ is Morse).  Then the Poincar\'e--Hopf theorem
applied to the vector field $\nabla^h s$ gives
\[
\sum_{x\in\Crit(s)} \operatorname{ind}_x(\nabla^h s) = \chi(M).
\]
In particular, since $\Crit(s) \subseteq \mathcal{D}(\Phi,s)$ by
Theorem~\ref{thm:morse-degen}, the degeneration locus contains a finite set
of points whose total index equals $\chi(M)$.  If $s$ is Morse, then each
index is $\pm 1$ and the identity reduces to the alternating-sum formula for
the Euler characteristic.
\end{corollary}

For the $n$-torus $T^n$, $\sum b_k = 2^n$, so any gradient-dominated
rule requires at least $2^n$ degeneration points for a Morse seam.
For $S^2$, $\sum b_k = 2$: the two critical points of any
height function are the irreducible degeneration locus of the
gradient metric---precisely the poles where the score geometry
becomes singular.

\subsection{The Hessian degeneration locus and optimal transport}
\label{sec:OT}

For the Hessian rule ($\gamma = 1$, $\alpha = \beta = 0$), the
degeneration locus $\mathcal{D} = \{\det(\nabla^2 s) = 0\}$ admits a
precise characterisation via optimal transport theory, which gives it
far more analytical substance than the abstract definition alone.

\begin{proposition}[Seam--optimal-transport correspondence]
\label{prop:OT}
Let $s\colon \mathbb{R}^n \to \mathbb{R}$ be a smooth strictly convex
function and $g = \nabla^2 s$ the associated Hessian metric.
\begin{enumerate}[label=(\roman*)]
\item The volume form of $g$ satisfies
  $dV_g = \sqrt{\det(\nabla^2 s)}\;dV_h$.  In particular, given
  a source measure $\mu_0 = \rho_0\,dV_h$ and a target measure
  $\mu_1 = \rho_1\,dV_h$, the condition $(\nabla s)_*\mu_0 = \mu_1$
  is equivalent to the \emph{Monge--Amp\`ere equation}
  \begin{equation}
  \label{eq:monge-ampere}
  \det(\nabla^2 s(x)) = \frac{\rho_0(x)}{\rho_1(\nabla s(x))}.
  \end{equation}
  By Brenier's theorem~\cite{Villani2009}, the unique $L^2$-optimal
  transport map
  $T = \nabla s$ between $\mu_0$ and $\mu_1$ arises from a convex
  potential~$s$---i.e., from a seam whose Hessian rule produces a
  positive-definite metric.

\item The degeneration locus of the Hessian rule is
  $\mathcal{D}(\Phi_{\mathrm{Hess}}, s) = \{\det(\nabla^2 s) = 0\}$,
  which in the optimal-transport setting is exactly the
  \emph{singular support} of the transport plan: the set where the
  transport map $\nabla s$ fails to be a local diffeomorphism.
\end{enumerate}
\end{proposition}

\begin{proof}
For~(i): in coordinates, $\det g = \det(\nabla^2 s)$ and
$dV_g = \sqrt{\det g}\,dx^1 \wedge \cdots \wedge dx^n
= \sqrt{\det(\nabla^2 s)}\,dV_h$.
The change-of-variables formula for $T = \nabla s$ gives
$\int_A \rho_1(y)\,dy = \int_{T^{-1}(A)} \rho_0(x)\,dx$,
which is~\eqref{eq:monge-ampere}.

For~(ii): the Jacobian of $T = \nabla s$ is $\nabla^2 s$; the map
fails to be a local diffeomorphism exactly where this Jacobian is
singular, i.e.\ $\det(\nabla^2 s) = 0$.
\end{proof}

\begin{remark}[Displacement interpolation as affine seams in Euclidean space]
On $\mathbb{R}^n$ with the quadratic cost, McCann's displacement
interpolation between the identity map and an optimal transport map
$T=\nabla\varphi$ is
$T_t(x)=(1-t)x+t\nabla\varphi(x)$.
Since $T_t=\nabla s_t$ with
$s_t(x)=(1-t)\tfrac12|x|^2+t\varphi(x)$,
the displacement geodesic corresponds to an \emph{affine} path in the
space of convex seams.  Under the Hessian rule this induces the linear
interpolation of Hessian metrics
$\nabla^2 s_t = (1-t)\,I + t\,\nabla^2\varphi$.
\end{remark}

This identification imports the extensive regularity theory for
Monge--Amp\`ere equations into the seam-rule framework:

\begin{corollary}[Regularity of the Hessian degeneration locus]
\label{cor:caffarelli}
Under the hypotheses of Proposition~\ref{prop:OT}, if $\rho_0$ and
$\rho_1$ are bounded away from zero and infinity on their supports,
then by Caffarelli's regularity theory~\cite{Caffarelli1992} for the
Monge--Amp\`ere equation, the Brenier potential $s$ is of class
$C^{1,\alpha}$ globally
and $C^\infty$ away from the boundary of the support.
In particular, the degeneration locus
$\mathcal{D} = \{\det(\nabla^2 s) = 0\}$ can occur only at the
\emph{boundary of the convex support} of the measures, and the Hessian
metric $g = \nabla^2 s$ is smooth in the interior.
\end{corollary}

\begin{remark}[Consequences for the framework]
The correspondence clarifies several aspects of the paper's machinery
in the Hessian case:
\begin{itemize}
\item The positive-definiteness condition $\nabla^2 s > 0$ for the
  Hessian rule is exactly Brenier's convexity condition on the
  optimal transport potential.
\item The inverse metric-design algorithm
  (Remark~\ref{rem:inverse-algorithm}) specialises, for the Hessian
  rule, to an iterative Monge--Amp\`ere solver; the normal-equation
  approach connects to the known numerical methods for optimal
  transport (Benamou--Brenier, entropic regularisation, etc.).
\item The information-geometry example
  (\S\ref{sec:ex-fisher}) is a special case: the log-partition
  function $A(\theta)$ of an exponential family is a convex potential,
  and the Fisher--Rao metric $\nabla^2 A$ is the Hessian metric of
  the trivial transport on the parameter space.
\end{itemize}
\end{remark}

% ================================================================
\section{Composition and Algebra of Rules}
\label{sec:composition}
% ================================================================

Rules can be combined in several ways, each producing a new rule.

\begin{definition}[Sequential composition]
If $\Phi_1$ produces a background suitable for $\Phi_2$, the
composition
$\Phi_{2 \circ 1}(s_1, s_2;\, h)
:= \Phi_2(s_2;\, \Phi_1(s_1;\, h))$
is a rule with two seam inputs.  The linearisation of $\Phi_{2\circ 1}$
at $(s_1, s_2)$ involves the chain rule:
$D\Phi_{2\circ 1} = D_2 \Phi_2 \circ D_1 \Phi_1
+ D_{\text{bg}} \Phi_2 \circ D\Phi_1$.
\end{definition}

\begin{definition}[Linear combination]
Given rules $\Phi^{(1)}, \ldots, \Phi^{(m)}$ all targeting
$\Sym^2(T^*M)$, a \emph{mixed rule} is
$\Phi_{\text{mix}} = \sum_{i=1}^m c_i \Phi^{(i)}$
with $c_i \in C^\infty(M)$ (or constants).  By the Universal
Decomposition, every order-$\le 2$ isotropic rule is already a mixed
rule with $m = 3$.
\end{definition}

\begin{example}[The Ricci flow as sequential composition]
The (unnormalised) Ricci flow
$\partial_t g = -2\Ric(g)$ can be viewed, at each instant, as a
second-order rule on the local potential representation
$g = \nabla^2 s$: the Ricci tensor of a Hessian metric depends on
the third jet $\nabla^3 s$ (cf.\ Proposition~\ref{prop:curv-general}).
Writing the flow directly as an evolution equation for $s$ requires a
further reduction (e.g. applying a normal operator), after which one
obtains a fourth-order parabolic scalar seam equation.
\end{example}

\begin{proposition}[Composition closure and jet order]
\label{prop:composition-closure}
Sequential composition of rules interacts with the jet filtration
as follows.
\begin{enumerate}[label=\textup{(\alph*)}]
\item \textbf{Conformal--conformal.}
  $e^{2s_2}(e^{2s_1}h) = e^{2(s_1+s_2)}h$; the composed rule
  remains conformal with composed seam $s_1 + s_2$.  The
  conformal rule is closed under sequential composition.
\item \textbf{Hessian--Hessian.}
  If $g_1 = \nabla^2_h s_1$ and $g_2 = \nabla^2_{g_1} s_2$,
  then
  \[
  g_2 = \nabla^2_h s_2
    - \Gamma(g_1,h)^k{}_{ij}\,\nabla_k s_2,
  \]
  where $\Gamma(g_1,h)$ is the connection difference tensor,
  which involves $\nabla^3_h s_1$ through the Christoffel
  symbols of~$g_1$.  Thus the composed metric escapes to
  the \emph{order-3} family
  (Proposition~\ref{prop:order3}), specifically involving
  the generator $\Pi_{\nabla^3,2}$.
  The Hessian rule is \emph{not} closed under composition.
\item \textbf{Conformal--Hessian.}
  The conformal change of connection gives
  \[
  \nabla^2_{e^{2s_1}h} s_2
  = \nabla^2_h s_2
    - 2\,\langle ds_1, ds_2\rangle_h\, h
    + 2\, ds_1 \odot ds_2,
  \]
  which is a linear combination of $h$, $ds_1 \odot ds_2$,
  and $\nabla^2_h s_2$---all order-$\le 2$ tensors.
  Conformal--Hessian composition remains within the
  three-term classification.
\end{enumerate}
\end{proposition}

\begin{proof}[Proof sketch]
Part~(a) is immediate.  For~(b), $\nabla^{g_1} = \nabla^h +
\Gamma(g_1,h)$ and the connection difference tensor
$\Gamma(g_1,h)^k{}_{ij} = \tfrac{1}{2}(g_1)^{kl}
(\nabla_i(g_1)_{jl} + \nabla_j(g_1)_{il} - \nabla_l(g_1)_{ij})$
involves $\nabla^3_h s_1$ via the derivatives of
$(g_1)_{ij} = (\nabla^2_h s_1)_{ij}$.
For~(c), the classical conformal Christoffel formula
$\Gamma(e^{2s_1}h,h)^k{}_{ij}
= \delta^k_i\,\partial_j s_1 + \delta^k_j\,\partial_i s_1
- h_{ij}\,h^{kl}\partial_l s_1$
contains only first derivatives of~$s_1$; substituting into
$\nabla^2_{g_1} s_2 = \nabla^2_h s_2
- \Gamma(g_1,h)^k{}_{ij}\partial_k s_2$ gives the stated formula.
\end{proof}

Parts (b) and~(c) reveal an asymmetry: repeated Hessian composition
``climbs the jet hierarchy,'' escaping the order-2 family, while
conformal--Hessian mixing is stable at order~2.  This gives concrete
guidance for multi-stage seam constructions: interleaving conformal
steps preserves tractability, whereas chaining Hessian steps
necessarily introduces higher-order jets.

% ================================================================
\section{The Seam Fibre}
\label{sec:fibre}
% ================================================================

The Universal Decomposition shows that a single seam generates a
whole family of geometries as the rule varies.  We now formalise this
observation.

\begin{definition}[Seam fibre]
\label{def:fibre}
Let $(M,h)$ be a Riemannian background and $s \in C^\infty(M)$ a
seam.  The \emph{seam fibre} of~$s$ is the set of all symmetric
2-tensors obtainable from~$s$ by admissible rules of order~$\le 2$:
\[
\mathcal{F}(s;\,h)
:= \bigl\{\,\alpha\,h + \beta\,ds\otimes ds + \gamma\,\nabla^2 s
   \;:\; (\alpha,\beta,\gamma) \in
   C^\infty(\mathcal{I})^3\bigr\},
\]
where $\mathcal{I} = (s,\,|\nabla s|_h^2,\,\Delta_h s)$ are the
$\OO(n)$-invariant jet scalars of~$s$.
\end{definition}

At each point $x \in M$, the fibre is the affine span of the three
tensors $h_x$, $(ds\otimes ds)_x$, and $(\nabla^2 s)_x$ with
coefficents depending smoothly on the jet invariants.  We now
identify which features of the geometry are intrinsic to the seam
and shared across the entire fibre, and which depend on the rule.

\begin{proposition}[Rule-invariant vs.\ rule-dependent features]
\label{prop:fibre-invariants}
Let $g_1 = \Phi_1(s;h)$ and $g_2 = \Phi_2(s;h)$ be two
Riemannian metrics in the same fibre $\mathcal{F}(s;h)$.  Then:
\begin{enumerate}[label=\textup{(\roman*)}]
\item \textbf{Shared (rule-invariant).}  $g_1$ and $g_2$ share:
  \begin{itemize}
  \item the level sets $\{s = c\}$ and their nesting;
  \item the critical set $\operatorname{Crit}(s) = \{\nabla s = 0\}$;
  \item the symmetry group
    $\operatorname{Isom}(M,h) \cap \operatorname{Sym}(s)$,
    where $\operatorname{Sym}(s) = \{\varphi : \varphi^* s = s\}$;
  \item the topological type of the sublevel sets $\{s \le c\}$.
  \end{itemize}
\item \textbf{Rule-dependent.}  The following can differ between
  $g_1$ and $g_2$:
  \begin{itemize}
  \item Riemannian curvature, geodesics, and distances;
  \item metric signature (e.g.\ the Hessian rule may give
    Lorentzian signature while the conformal rule gives
    Riemannian);
  \item the degeneration locus $\mathcal{D}(\Phi_i, s)$
    (e.g.\ the gradient rule degenerates at
    $\operatorname{Crit}(s)$, while the conformal rule
    $e^{2s}h$ is everywhere nondegenerate).
  \end{itemize}
\end{enumerate}
\end{proposition}

\begin{proof}
Part~(i): the level sets, critical set, and sublevel-set topology
depend only on~$s$ (a smooth function), not on the rule.  If
$\varphi$ is an isometry of~$h$ with $\varphi^* s = s$, then by
naturality (Axiom~\ref{item:naturality}),
$\Phi(\varphi^* s;\,\varphi^* h) = \varphi^*\Phi(s;h)$ for any
rule~$\Phi$, so $\varphi$ is an isometry of~$g$ as well.

Part~(ii): the conformal rule with $\alpha = e^{2s}$ gives
positive-definite $g_1 = e^{2s}h$ everywhere.  The Hessian rule
gives $g_2 = \nabla^2 s$, which is indefinite whenever $\nabla^2 s$
has mixed eigenvalues (e.g.\ for harmonic~$s$, see
\S\ref{sec:ex-hessian}).  The degeneration loci differ because
each rule's symbol has a different kernel structure
(Proposition~\ref{prop:symbol}).
\end{proof}

\begin{remark}[The seam as a latent geometric reservoir]
Conceptually, the 2-jet $j^2 s$ at each point simultaneously encodes
three independent tensorial objects---$s$ (a scalar), $ds \otimes ds$
(a rank-1 tensor), and $\nabla^2 s$ (a full symmetric tensor).
A rule is a choice of which of these to ``listen to'' and with what
weight.  The seam therefore does not generate one geometry but an
entire family, parameterised by the coefficient functions
$(\alpha,\beta,\gamma)$.  Different rules extract different aspects
of the same underlying scalar data---an observation that is
practically useful whenever one rule produces better-conditioned
geometry than another for a given computational task
(see Remark~\ref{rem:inverse-algorithm}).
\end{remark}

The seam fibre raises a natural uniqueness question: if two seams
produce the same geometry under the same rule, how are they related?

\begin{theorem}[Rigidity of seam representations]
\label{thm:rigidity}
Let $(M,h)$ be a connected compact manifold without boundary and
let $\Phi = \alpha\,h + \beta\,ds\otimes ds + \gamma\,\nabla^2 s$
be a rule with coefficients depending smoothly on
$(s,|\nabla s|^2)$.  Suppose $\Phi(s_1;h) = \Phi(s_2;h)$
pointwise on~$M$.
\begin{enumerate}[label=\textup{(\alph*)}]
\item \textbf{Conformal rule} ($\beta = \gamma = 0$,
  $\alpha = e^{2s}$):
  $s_1 - s_2$ is constant on~$M$.  The gauge group is
  $\mathbb{R}$ (additive translations).
\item \textbf{Hessian rule} ($\alpha = \beta = 0$, $\gamma = 1$):
  $s_1 - s_2$ is an affine function
  ($\nabla^2(s_1 - s_2) = 0$), so
  $\dim\ker = n + 1$ on $\mathbb{R}^n$ and
  $\dim\ker = 1$ (constants) on any connected compact manifold
  without boundary.
\item \textbf{General rule with $\gamma \neq 0$}:
  $\delta s := s_1 - s_2$ satisfies
  $D\Phi_{s_1}(\delta s) = O(|\delta s|^2)$.
  In particular, the space of infinitesimal gauge
  symmetries is $\ker D\Phi_{s_1}$, which is
  finite-dimensional by the Fredholm property
  (Proposition~\ref{prop:fredholm}).
  Generically $\dim\ker D\Phi_{s_0} = 0$ or~$1$, so the
  seam representation is \emph{rigid}: at most a
  one-parameter family of seams produces the same geometry.
\end{enumerate}
\end{theorem}

\begin{proof}
Subtract: $\Phi(s_1;h) - \Phi(s_2;h) = 0$.
Expanding about $s_1$,
$D\Phi_{s_1}(s_2 - s_1) + Q(s_1, s_2) = 0$
where $Q$ is the quadratic (and higher) remainder.

For~(a): $e^{2s_1} = e^{2s_2}$ on connected~$M$ forces
$s_1 = s_2 + c$ for a constant~$c \in \mathbb{R}$.

For~(b): $\nabla^2 s_1 = \nabla^2 s_2$ gives
$\nabla^2(s_1 - s_2) = 0$, so $s_1 - s_2$ is affine.
Taking the trace yields $\Delta_h(s_1-s_2)=0$.
On a connected compact manifold without boundary, harmonic functions
are constant, hence $s_1-s_2$ is constant.

For~(c): the linearised equation
$D\Phi_{s_1}(\delta s) = -Q(s_1, s_2)$ has $Q = O(|\delta s|^2)$.
For infinitesimal gauges ($Q = 0$), the solution space is
$\ker D\Phi_{s_1}$.  The normal operator
$(D\Phi_{s_1})^* D\Phi_{s_1}$ is fourth-order elliptic with
principal symbol $\gamma^2|\xi|^4$
(Proposition~\ref{prop:fredholm}), hence Fredholm of index~$0$.
Self-adjointness gives
$\dim\ker = \dim\ker N_{s_1} = \dim\coker N_{s_1}$,
and for generic~$s_1$ the kernel is trivial or one-dimensional
(corresponding to the residual constant gauge).
\end{proof}

% ================================================================
\section{Universality: When Can a Rule Reach All Geometries?}
\label{sec:universality}
% ================================================================

\begin{definition}[Universality]
A rule $\Phi$ on $(M,h)$ is \emph{universal (up to diffeomorphism)}
if for every Riemannian metric $g$ on $M$ there exist a seam $s$
and a diffeomorphism $\varphi$ such that
$g = \varphi^* \Phi(s;\, h)$.
\end{definition}

\begin{proposition}[No single order-$\le 2$ basis rule is universal in
  dimension $n \ge 3$]
\label{prop:no-universal}
On a compact manifold $M^n$ with $n \ge 3$:
\begin{enumerate}[label=(\alph*)]
\item The conformal rule $\Phi_\text{Conf}(s;h) = e^{2s}h$ produces
  only metrics conformal to $h$.  The space of conformal classes on a
  manifold with $n \ge 3$ is infinite-dimensional, but each class is
  parametrised by a single function.  Hence
  $\image(\Phi_\text{Conf}) = [h]$, a single conformal class.
\item The Hessian rule $\Phi_\text{Hess}(s) = \nabla^2 s$ produces
  only exact Hessians.  Such tensors satisfy overdetermined
  compatibility identities (generalized Saint-Venant conditions):
  on a flat background they reduce to
  $\partial_i \partial_j g_{kl} = \partial_k \partial_l g_{ij}$,
  while on a curved background the corresponding identities involve
  the curvature $\Rm(h)$ through commutators of covariant derivatives.
  Generically, metrics do not satisfy these.
\item The gradient rule $\Phi_\text{Grad}(s;h) = |\nabla s|^2 h$
  lies within the conformal class of $h$, and the conformal factor
  $|\nabla s|^2$ vanishes at every critical point of $s$.  Hence it
  cannot produce a metric with everywhere-positive conformal factor
  unless $s$ has no critical points---impossible on a compact manifold
  (by Ljusternik--Schnirelmann theory, $s$ has at least $\text{cat}(M)
  \ge 2$ critical points).
\end{enumerate}
\end{proposition}

\begin{proposition}[Conformal universality on genus-$0$ surfaces]
\label{prop:genus0}
On a closed surface of genus $0$, the conformal rule is universal up
to diffeomorphism: every Riemannian metric $g$ can be written as
$g = e^{2s}\varphi^* h_0$ for some seam $s$ and diffeomorphism
$\varphi$.  This is a consequence of
the uniformisation theorem and the rigidity of the conformal structure
on $S^2$.
\end{proposition}

We now give obstructions to universality for order-$\le 2$ seam rules.
On surfaces, the extended linearised system is \emph{never elliptic}
when a Hessian term is present ($\gamma\neq 0$): its
Douglis--Nirenberg principal symbol has a non-trivial kernel for every
$\xi\neq 0$.  In higher dimensions, a \emph{jet-dimension deficit}
blocks density for any fixed rule.

\begin{theorem}[No-go: Douglis--Nirenberg symbol degeneracy on surfaces]
\label{thm:no-go-surface}
Let $(M^2, h)$ be a closed oriented surface.
Consider the extended seam map
\[
F(s,\varphi) := \varphi^*\bigl[\alpha\,h + \beta\,ds\otimes ds + \gamma\,\nabla^2 s\bigr],
\qquad \gamma \neq 0,
\]
and its linearisation at a background seam~$s_0$ and the identity
diffeomorphism:
\[
DF_{(s_0,\id)}(\dot s, X)
= \alpha'\dot s\, h + \beta'\dot s\,ds_0\otimes ds_0
  + 2\beta\,ds_0\odot d\dot s + \gamma\,\nabla^2\dot s + \mathcal{L}_X\,\Phi(s_0;h).
\]
Then the Douglis--Nirenberg principal symbol
\[
\sigma(\xi)(\dot s, X)
= \gamma\,(\xi \otimes \xi)\,\dot s
  + \xi \otimes X^\flat + X^\flat \otimes \xi
\]
has a non-trivial kernel for every $\xi\neq 0$, spanned by
$(\dot s,\, -\tfrac{\gamma}{2}\,\xi\,\dot s)$.  In particular,
the symbol is not surjective for any $\xi\neq 0$.
\end{theorem}

\begin{proof}
The principal symbol of the Lie-derivative map
$X \mapsto \mathcal{L}_X h$ at covector $\xi$ is
$\xi \otimes X^\flat + X^\flat \otimes \xi$.
Hence the Douglis--Nirenberg principal symbol of
$(\dot s,X) \mapsto \gamma\nabla^2\dot s + \mathcal{L}_X h$
is
$\gamma\,(\xi\otimes\xi)\,\dot s + \xi\otimes X^\flat + X^\flat\otimes\xi$.
Choosing $X^\flat = -\tfrac{\gamma}{2}\,\xi\,\dot s$ cancels the
Hessian contribution:
\[
\gamma\,\xi_i\xi_j\,\dot s
- \tfrac{\gamma}{2}\,\xi_i\xi_j\,\dot s
- \tfrac{\gamma}{2}\,\xi_j\xi_i\,\dot s = 0.
\]
Since the symbol maps
$\mathbb{R}^1 \times \mathbb{R}^2 \to \Sym^2(\mathbb{R}^2)$
(both 3-dimensional), the 1-dimensional kernel implies the
image has dimension~2: the symbol is not surjective.
\end{proof}

\begin{remark}[What the symbol degeneracy does and does not say]
Theorem~\ref{thm:no-go-surface} is a \emph{non-ellipticity} statement:
it rules out applying standard elliptic Fredholm/implicit-function
theory to the extended surface system when $\gamma\neq 0$.
It does \emph{not} by itself preclude local surjectivity, since
lower-order terms may compensate for the symbol degeneracy.
Indeed, on the torus the gradient-tensor term ($\beta\neq 0$)
produces Teichm\"uller directions and yields local universality
(Theorem~\ref{thm:torus} below).
\end{remark}

This dichotomy motivates separating two questions on surfaces: the
principal-symbol analysis (which is obstructed whenever $\gamma\neq 0$)
and the finite-dimensional Teichm\"uller directions (which may still be
accessible through the lower-order gradient-tensor term when $\beta\neq 0$).
Theorem~\ref{thm:torus} gives a clean model case where this compensation
happens and the full linearisation becomes surjective.  For genus
$g\ge 2$, whether the $\beta$-term can generically navigate the
$(6g{-}6)$ Teichm\"uller directions remains open.

\begin{theorem}[Local universality on the torus]
\label{thm:torus}
Let $(\mathbb{T}^2, h)$ be a flat torus and fix a rule
with $\alpha'(s_0) \neq 0$ and $\beta(s_0) \neq 0$ for
some background seam~$s_0$.  If the trace-free Hessian
$(\nabla^2 s_0)^{TF}$ is not everywhere proportional to
a single constant traceless tensor\textup{---}in particular,
if $s_0$ is any Morse function with at least two
saddle points of distinct orientation\textup{---}then the
extended seam map
\[
F(s,\varphi) = \varphi^*\bigl[\alpha\,h + \beta\,ds\otimes ds
  + \gamma\,\nabla^2 s\bigr]
\]
is a local submersion at $(s_0, \id)$: its image contains
a $C^k$-open neighbourhood of\/
$g_0 = \Phi(s_0;h)$ in\/ $\Met(\mathbb{T}^2)$.

In particular, the three-term family is \emph{locally
universal on the torus}.
\end{theorem}

\begin{proof}
The tangent space $T_{g_0}\Met(\mathbb{T}^2)$ decomposes
into three summands with respect to the $L^2(g_0)$-Hodge
splitting:
\begin{enumerate}[label=(\roman*)]
\item \emph{Diffeomorphism directions:}
  $\image(\delta^*_{g_0})$, where
  $\delta^*_{g_0}(X) = \tfrac{1}{2}\mathcal{L}_X g_0$.
\item \emph{Conformal directions:}
  $\{f\, g_0 : f \in C^\infty(\mathbb{T}^2)\}$
  (within $\ker\delta_{g_0}$).
\item \emph{Teichm\"uller directions:}
  The space $H_{TT}(h)$ of constant traceless symmetric
  2-tensors with respect to~$h$, a real vector space of
  dimension~2, spanned by
  $e_1 = dx^2 - dy^2$ and $e_2 = 2\,dx\,dy$.
\end{enumerate}
We show the linearisation
$DF_{(s_0,\id)}\colon C^\infty(\mathbb{T}^2) \times
\mathfrak{X}(\mathbb{T}^2) \to \Gamma(\Sym^2 T^*\mathbb{T}^2)$
surjects onto each summand.

\medskip\noindent
\emph{Summand~(i).}
The variation with respect to a vector field~$X$ gives
$DF(0,X) = \mathcal{L}_X g_0$, which spans
$\image(\delta^*_{g_0})$ by definition.

\medskip\noindent
\emph{Summand~(ii).}
The seam variation is
\begin{equation}
\label{eq:torus-lin}
D_s\Phi_{s_0}(\dot s)
= \alpha'\dot s\;h
  + \beta'\dot s\;ds_0\otimes ds_0
  + 2\beta\;ds_0 \odot d\dot s
  + \gamma\;\nabla^2\dot s,
\end{equation}
where $\alpha' = \partial\alpha/\partial s$ evaluated at
the invariants of~$s_0$, and similarly for $\beta'$.
The first two terms are pointwise multiples of symmetric
tensors, so taking
$\dot s = f/(\alpha' + \beta'|\nabla s_0|^2)$ for any
$f \in C^\infty$ yields a trace contribution
$f\cdot(h + \tfrac{\beta'}{\alpha'+\beta'|\nabla s_0|^2}
\,ds_0{\otimes}ds_0)$ plus derivative terms.
Since $\alpha'(s_0) \neq 0$, the trace
part spans an open cone in the conformal directions; the
derivative terms lie in summands~(i) and~(iii) and are
handled separately.

\medskip\noindent
\emph{Summand~(iii): the key step.}
We project $D_s\Phi_{s_0}(\dot s)$ onto $H_{TT}(h)$.
Because $e_a \in H_{TT}$ is constant and traceless, the
inner product with the first two terms of~\eqref{eq:torus-lin}
vanishes: $\alpha'\dot s\,\langle h, e_a\rangle = 0$
(trace vs.\ traceless), and
$\beta'\dot s\,\langle ds_0{\otimes}ds_0, e_a\rangle$
has the same status once we integrate the zero-order part
(it contributes to the component we control via~$\dot s$
below).  The $\gamma\nabla^2\dot s$ term projects to zero
in $H_{TT}$ because on the flat torus
$\langle \nabla^2\dot s, e_a\rangle_{L^2}
= \int e_a^{ij}\nabla_i\nabla_j\dot s\,dA
= \int \dot s\,\nabla_j\nabla_i e_a^{ij}\,dA = 0$
(since $e_a$ is constant).

The sole non-trivial TT contribution comes from the
gradient-tensor variation $2\beta\,ds_0\odot d\dot s$:
\begin{align*}
\langle 2\beta\,ds_0\odot d\dot s,\; e_a\rangle_{L^2}
&= 2\beta \int_{\mathbb{T}^2} e_a^{ij}\,
   \nabla_i s_0\,\nabla_j \dot s\;dA \\
&= -2\beta \int_{\mathbb{T}^2} \dot s\;
   e_a^{ij}\,\nabla_i\nabla_j s_0\;dA
   \qquad\text{(integration by parts; $e_a$ constant)}\\
&= -2\beta \int_{\mathbb{T}^2} \dot s\;
   \bigl[e_a : (\nabla^2 s_0)^{TF}\bigr]\;dA,
\end{align*}
where $e_a : (\nabla^2 s_0)^{TF}
= e_a^{ij}(\nabla^2 s_0)^{TF}_{ij}$
(the trace part drops out since $e_a$ is traceless).
Define $\lambda_a := e_a : (\nabla^2 s_0)^{TF}
\in C^\infty(\mathbb{T}^2)$.
The map
$\dot s \mapsto
\bigl(\langle\cdot,e_1\rangle,\;\langle\cdot,e_2\rangle\bigr)
\in \mathbb{R}^2$
is then $\dot s \mapsto
-2\beta\bigl(\int\dot s\,\lambda_1\,dA,\;
\int\dot s\,\lambda_2\,dA\bigr)$.

This map surjects onto $\mathbb{R}^2$ if and only if
$\lambda_1$ and $\lambda_2$ are \emph{linearly independent
as functions} on $\mathbb{T}^2$.  We verify this for the
explicit seam
$s_0(x,y) = \cos(2\pi x) + \cos(2\pi(x+y))$:
\begin{align*}
\lambda_1
&= (\partial_{xx}s_0 - \partial_{yy}s_0)
= -4\pi^2\cos(2\pi x), \\[2pt]
\lambda_2
&= 2\,\partial_{xy}s_0
= -8\pi^2\cos\bigl(2\pi(x+y)\bigr).
\end{align*}
These are distinct Fourier modes on $\mathbb{T}^2$,
hence linearly independent in $L^2$.  Therefore
the TT-projection of the image of
$D_s\Phi_{s_0}$ spans all of $H_{TT}(h)$.

\medskip\noindent
\emph{Application of the inverse function theorem.}
The three summands are covered: (i)~by diffeomorphisms,
(ii)~by the conformal part of the seam variation, and
(iii)~by the gradient-tensor part.  The linearised map
$DF_{(s_0,\id)}$ is therefore surjective.  A bounded
right inverse exists because the diffeomorphism and
conformal parts are handled by standard elliptic theory
(the vector Laplacian and scalar identity, respectively),
while the TT part is finite-dimensional.  The
Nash--Moser inverse function theorem on tame Fr\'echet
spaces gives local surjectivity of the nonlinear map.
For the extended map incorporating diffeomorphisms, one cannot
in general replace Nash--Moser by a standard Banach-space implicit
function theorem on fixed Sobolev spaces without an additional
gauge-fixing (e.g.\ of DeTurck type), since the pullback action
$\varphi \mapsto \varphi^*g$ loses derivatives.
\end{proof}

\begin{remark}[The gradient tensor as Teichm\"uller navigator]
\label{rem:beta-hero}
Theorem~\ref{thm:torus} identifies a precise role for the
gradient-tensor direction $ds\otimes ds$, noted in
Remark~\ref{rem:what-classif} as ``largely unexplored.''
On the torus, the conformal rule changes the conformal factor
but cannot shift the conformal class; the Hessian rule is
gauge-equivalent to a diffeomorphism at highest order.
It is the $\beta$-term alone
that produces traceless deformations capable of navigating
across Teichm\"uller space.  The genericity condition on~$s_0$
(that its trace-free Hessian sweeps both TT directions) is
mild: any Morse function whose gradient rotates through at
least two linearly independent directions suffices.
The condition fails only on a set of infinite codimension
in $C^\infty(\mathbb{T}^2)$.
\end{remark}

\begin{proposition}[A seam with independent TT pairings on a hyperbolic surface]
\label{prop:hyperbolic-tt-independence}
Let $(M_g,h)$ be a closed hyperbolic surface of genus $g\ge 2$.
Let $\{e_a\}_{a=1}^{6g-6}$ be any $L^2(h)$-orthonormal basis of the
transverse--traceless space
$H_{TT}(h)=\ker(\delta_h)\cap\ker(\tr_h)$.
Then there exists a seam $s_0\in C^\infty(M_g)$ such that the functions
\[
\lambda_a:= e_a^{ij}\nabla_i\nabla_j s_0
= e_a:(\nabla_h^2 s_0)^{TF}
\in C^\infty(M_g),\qquad a=1,\ldots,6g-6,
\]
are linearly independent in $L^2(M_g,dA_h)$.
\end{proposition}

\begin{proof}
Define the evaluation map on the unit tangent bundle
$S(TM_g):=\{(p,\hat\xi):p\in M_g,\ \hat\xi\in T_pM_g,\ |\hat\xi|_h=1\}$,
\[
F:S(TM_g)\to\mathbb{R}^{6g-6},\qquad
F(p,\hat\xi):=\bigl(e_a(p)(\hat\xi,\hat\xi)\bigr)_{a=1}^{6g-6}.
\]
We claim that $\mathrm{span}(\mathrm{im}(F))=\mathbb{R}^{6g-6}$.
Indeed, if $c=(c_a)\neq 0$ annihilates $\mathrm{im}(F)$, then the TT tensor
$\Omega:=\sum_a c_a e_a\in H_{TT}(h)$ satisfies $\Omega(p)(\hat\xi,\hat\xi)=0$
for every $(p,\hat\xi)\in S(TM_g)$.  On a surface, a traceless symmetric
bilinear form is determined by its values on unit directions, so this forces
$\Omega(p)=0$ for all $p$, hence $\Omega=0$, contradiction.
Therefore $\mathrm{im}(F)$ spans.

Choose $6g-6$ points $(p_b,\hat\xi_b)\in S(TM_g)$ such that the vectors
$F(p_b,\hat\xi_b)$ are linearly independent.  Equivalently, the matrix
\[
B_{ab}:=e_a(p_b)(\hat\xi_b,\hat\xi_b)
\]
is invertible.  Choose the $p_b$ distinct and supported in pairwise disjoint
small $h$-geodesic balls.

For each $b$, pick geodesic normal coordinates $x=(x^1,x^2)$ centred at $p_b$.
Let
$q_b(x):=\tfrac12(x\cdot\hat\xi_b)^2$ and choose a cutoff
$\chi_b\in C_c^\infty$ supported in that ball with $\chi_b(p_b)=1$ and
$d\chi_b(p_b)=0$.  Set $s_b:=\chi_b q_b$.  Since $\Gamma^k_{ij}(p_b)=0$ in
normal coordinates and $d\chi_b(p_b)=0$, we have
$(\nabla^2 s_b)(p_b)=\hat\xi_b\otimes\hat\xi_b$.  Because each $e_a(p_b)$ is
traceless, contracting kills the trace part, so
\[
\lambda_a(s_b)(p_b)=e_a^{ij}(p_b)(\nabla_i\nabla_j s_b)(p_b)=B_{ab}.
\]
Let $s_0:=\sum_{b=1}^{6g-6}s_b$.  By disjointness of supports,
$\lambda_a(s_0)(p_b)=B_{ab}$ for every $a,b$.

If $\sum_a c_a\lambda_a(s_0)=0$ in $L^2$, then the smooth function
$\sum_a c_a\lambda_a(s_0)$ vanishes almost everywhere and hence everywhere.
Evaluating at $p_b$ yields $\sum_a c_a B_{ab}=0$ for all $b$, i.e.
$B^T c=0$.  Since $B$ is invertible, $c=0$.  Thus the $\lambda_a$ are
$L^2$-independent.
\end{proof}

\begin{remark}[A conditional route to higher-genus local universality]
Proposition~\ref{prop:hyperbolic-tt-independence} supplies the key
finite-dimensional spanning statement at the \emph{reference} hyperbolic
metric $h$.  To turn this into a genus-$g\ge 2$ analogue of
Theorem~\ref{thm:torus}, one must transfer TT-surjectivity from $H_{TT}(h)$
to the actual base metric $g_0=\Phi(s_0;h)$.

The natural analytic mechanism is to view
\[
H_{TT}(g)=\ker(\tr_g)\cap\ker(\delta_g)=\ker P_g,\qquad
P_g:=(\tr_g,\delta_g),
\]
and to use the self-adjoint elliptic operator $A_g:=P_g^*P_g$.
On a closed surface of genus $\ge 2$, $\dim H_{TT}(g)=6g-6$ for every metric
$g$, and there are no nontrivial conformal Killing fields.
The following stability statement is standard in elliptic perturbation theory.

\begin{lemma}[Stability of TT spectral projections]
\label{lem:tt-projection-stability}
Let $(M_g,h)$ be a closed hyperbolic surface of genus $g\ge 2$.
Fix an integer $k\ge 0$.
Then there exists an integer $m\ge k+3$ and a $C^m$-neighbourhood
$\mathcal{V}$ of $h$ in $\Met(M_g)$ such that for every $g\in\mathcal{V}$
there is a bounded projection
\[
\Pi^{TT}_g:H^k(\Sym^2T^*M_g)\to H^k(\Sym^2T^*M_g)
\]
with $\mathrm{ran}(\Pi^{TT}_g)=H_{TT}(g)$.
Moreover, $g\mapsto \Pi^{TT}_g$ is continuous in operator norm on
$\mathcal{L}(H^k,H^k)$.

In particular, for $g$ sufficiently close to $h$, the restricted map
$\Pi^{TT}_g\big|_{H_{TT}(h)}:H_{TT}(h)\to H_{TT}(g)$ is an isomorphism.
\end{lemma}

\begin{proof}[Proof sketch]
Set $P_g:=(\tr_g,\delta_g)$ and $A_g:=P_g^*P_g$.
Then $A_g$ is elliptic and self-adjoint on $L^2(g)$, and
$\ker(A_g)=\ker(P_g)=H_{TT}(g)$.
For $h$ fixed, let $\lambda_1(h)>0$ be the first positive eigenvalue of $A_h$.
Choose a small contour $\Gamma\subset\mathbb{C}$ enclosing $0$ but no other
point of $\mathrm{spec}(A_h)$ (for instance a circle of radius
$\lambda_1(h)/2$).
For $g$ sufficiently $C^m$-close to $h$, the coefficients of $A_g$ are close
to those of $A_h$ in the operator topology between Sobolev spaces, so the
resolvent $(A_g-z)^{-1}$ exists for $z\in\Gamma$ and is uniformly bounded on
$H^k$.  Define the spectral projection
\[
\Pi^{TT}_g:=\frac{1}{2\pi i}\oint_{\Gamma}(A_g-z)^{-1}\,dz.
\]
Standard resolvent identities then give continuity of
$g\mapsto\Pi^{TT}_g$ in operator norm.
The last claim follows because $\Pi^{TT}_g$ is a small perturbation of
$\Pi^{TT}_h$ and both ranges have the same finite dimension $6g-6$.
\end{proof}

Assuming such TT-projection stability, one can scale $s_0\mapsto\varepsilon s_0$
so that $g_\varepsilon=\Phi(\varepsilon s_0;h)$ is close to $h$, and then use
finite-dimensional perturbation stability to conclude that the TT component of
the linearised extended seam map at $g_\varepsilon$ remains surjective.
This would complete a higher-genus local universality theorem.
We leave this analytic stability step as part of the open surface universality
problem.
\end{remark}

For a \emph{fixed} rule in higher dimensions, a
complementary obstruction arises from dimension counting.

\begin{proposition}[Fixed-rule jet-dimension deficit in dimension $n \ge 3$]
\label{prop:jet-deficit}
Let $(M^n,h)$ be a compact manifold with $n \ge 3$ and let
$\Phi$ be any single rule (i.e.\ with fixed coefficient
functions $\alpha,\beta,\gamma$).  The extended seam map
$F(s,\varphi) = \varphi^*\Phi(s;h)$ maps $(n+1)$ scalar
functions to $\frac{n(n+1)}{2}$ independent metric components.
Since $\frac{n(n+1)}{2} > n+1$ for $n \ge 3$, the $r$-jet
spaces satisfy
\[
\dim J^r(\text{source})_x \sim (n+1)\binom{r+n}{n},
\qquad
\dim J^r(\text{target})_x \sim \tfrac{n(n+1)}{2}\binom{r+n}{n},
\]
with a strict deficit ratio
$\frac{n+1}{n(n+1)/2} = \frac{2}{n} < 1$ for $n \ge 3$.
Consequently, the image of a fixed rule has
\emph{infinite codimension} in $\Met(M)$: for any $r$,
the $r$-jet prolongation of~$F$ has rank strictly less than
the dimension of the target jet space, so the image cannot be
dense.
\end{proposition}

\begin{proof}
The jet-prolongation $j^r F$ maps
$J^r(C^\infty(M) \times \Diff(M))_x \to J^r(\Sym^2_+ T^*M)_x$.
The source has $n+1$ independent functions of $n$ variables,
contributing $(n+1)\binom{r+n}{n}$ free Taylor coefficients.
The target has $\frac{n(n+1)}{2}$ independent components, each
with $\binom{r+n}{n}$ Taylor coefficients.  For $n \ge 3$ and
all $r \ge 0$,
$(n+1) < \frac{n(n+1)}{2}$, so $j^r F$ cannot be surjective.
By the Thom transversality theorem (finite-dimensional version),
the image of $F$ is contained in a proper closed subset of
$C^k(\Sym^2_+ T^*M)$.
\end{proof}

Even allowing the coefficients to vary freely does not help
for $n \ge 4$.

\begin{theorem}[Non-universality with free coefficients, $n \ge 4$]
\label{thm:free-coeff-no-go}
Let $(M^n,h)$ be compact with $n \ge 4$.  Promote the
coefficients to \emph{completely free} scalar fields
$A,B,C \in C^\infty(M)$ (not required to depend on the jet
of~$s$) and consider the extended map
\[
  F(s,\varphi,A,B,C)
  = \varphi^*\bigl[A\,h + B\,ds{\otimes}ds + C\,\nabla^2 s\bigr].
\]
Even with this enlarged source, the image of~$F$ is
nowhere dense in $\Met(M)$.
\end{theorem}

\begin{proof}
Count independent functions of $n$ variables on the source and
target sides:
\[
  \underbrace{1}_{s} + \underbrace{n}_{\varphi}
  + \underbrace{3}_{A,B,C} = n + 4
  \qquad\text{vs.}\qquad
  \frac{n(n+1)}{2}.
\]
For $n \ge 4$, the inequality
$n + 4 < \frac{n(n+1)}{2}$
holds: the source has strictly fewer functional degrees of freedom.

At the $r$-jet level the prolongation
$j^r F \colon J^r(\text{source})_x \to J^r(\Sym^2_+ T^*M)_x$
maps between finite-dimensional spaces with
\[
  \dim J^r(\text{source})_x
  = (n{+}4)\tbinom{n{+}r}{n},
  \qquad
  \dim J^r(\text{target})_x
  = \tfrac{n(n{+}1)}{2}\tbinom{n{+}r}{n},
\]
and the source dimension is strictly smaller for every $r \ge 0$.
By the (finite-dimensional) Sard theorem, the image of the
jet prolongation has measure zero in the target jet space
for each~$r$.  Taking $r \to \infty$, the $C^\infty$-image
is contained in a meagre subset of $\Met(M)$.
\end{proof}

For $n = 3$ the dimension count reverses:
$n + 4 = 7 > 6 = n(n+1)/2$, so the system is formally
\emph{underdetermined} and no algebraic jet-dimension deficit
applies.  Nevertheless, a topological parity argument shows
that the Douglis--Nirenberg symbol is \emph{always} degenerate.

\begin{theorem}[Symbol parity obstruction in dimension $3$]
\label{thm:n3-parity}
Let $(M^3,h)$ be a closed $3$-manifold.  For any background
seam $s_0$ with $C_0 \neq 0$ and any free coefficient
functions $(A_0,B_0,C_0)$, the Douglis--Nirenberg principal
symbol of the linearised extended seam map
\[
  DF_{(s_0,\id,A_0,B_0,C_0)} \colon
  C^\infty(M) \times \mathfrak{X}(M) \times C^\infty(M)^3
  \;\longrightarrow\;
  \Gamma(\Sym^2 T^*M)
\]
fails to be surjective at some nonzero covector
$\xi_0 \in T^*_x M$ at every point~$x$.
Consequently, the system is never Douglis--Nirenberg elliptic.
\end{theorem}

\begin{proof}
Assign Douglis--Nirenberg source weights:
order~$2$ to~$\dot s$ (since $C_0\nabla^2\dot s$ is the
highest derivative), order~$1$ to the vector field~$X$
(from $\mathcal{L}_X g_0$), and order~$0$ to
$(\dot A,\dot B,\dot C)$.  With all target weights~$0$
the DN principal symbol at a nonzero covector $\xi$ is
\[
  \sigma(\xi)(\dot s,X,\dot A,\dot B,\dot C)
  = C_0\dot s\,(\xi{\otimes}\xi)
    + \xi \odot X
    + \dot A\,h
    + \dot B\,ds_0{\otimes}ds_0
    + \dot C\,\nabla^2 s_0\,.
\]

\medskip\noindent
\emph{Step~1: Absorption of the seam variable.}
Since $\xi{\otimes}\xi = \xi \odot \xi$, the seam
contribution $C_0\dot s\,(\xi{\otimes}\xi)$ lies in
the diffeomorphism orbit
$O_\xi = \{\xi \odot Y : Y \in \mathbb{R}^3\}$
(set $Y = C_0\dot s\,\xi$).  Defining
$\tilde X = X + C_0\dot s\,\xi$,
the symbol becomes
\[
  \sigma(\xi)
  = \xi \odot \tilde X
    + \dot A\,h
    + \dot B\,ds_0{\otimes}ds_0
    + \dot C\,\nabla^2 s_0\,.
\]
The seven-dimensional source
$(\dot s,X,\dot A,\dot B,\dot C) \in \mathbb{R}^7$
thereby collapses to the six-dimensional variable
$(\tilde X,\dot A,\dot B,\dot C) \in \mathbb{R}^6$;
the map $\sigma(\xi)$ factors through a
$6{\times}6$ matrix~$M(\xi)$.

\medskip\noindent
\emph{Step~2: Structure of $M(\xi)$.}
Work in normal coordinates at~$x$ with $h = \delta$.
Write $p = \nabla s_0 \in \mathbb{R}^3$,
$H = \nabla^2 s_0 \in \Sym^2(\mathbb{R}^3)$.
Index the six target components as
$(11),(22),(33),(12),(13),(23)$, and the source as
$(\tilde X_1,\tilde X_2,\tilde X_3,\dot A,\dot B,\dot C)$.
The explicit matrix is
\[
M(\xi) =
  \begin{pmatrix}
  \xi_1 & 0 & 0 & 1 & p_1^2 & H_{11}\\
  0 & \xi_2 & 0 & 1 & p_2^2 & H_{22}\\
  0 & 0 & \xi_3 & 1 & p_3^2 & H_{33}\\
  \frac{\xi_2}{2} & \frac{\xi_1}{2} & 0 & 0 & p_1p_2 & H_{12}\\
  \frac{\xi_3}{2} & 0 & \frac{\xi_1}{2} & 0 & p_1p_3 & H_{13}\\
  0 & \frac{\xi_3}{2} & \frac{\xi_2}{2} & 0 & p_2p_3 & H_{23}
  \end{pmatrix}.
\]
The first three columns are linear in~$\xi$; the last
three are independent of~$\xi$.

\medskip\noindent
\emph{Step~3: Homogeneity.}
Since columns $1$--$3$ scale linearly in~$\xi$ and
columns $4$--$6$ are constant,
$\det M(t\xi) = t^3\det M(\xi)$
for every $t \in \mathbb{R}$.
Thus $\det M$ is a \emph{homogeneous polynomial of
degree~$3$} in $(\xi_1,\xi_2,\xi_3)$.

\medskip\noindent
\emph{Step~4: Parity and the intermediate value theorem.}
Substituting $t = -1$:
\[
  \det M(-\xi) = (-1)^3\det M(\xi) = -\det M(\xi).
\]
That is, $f(\xi) := \det M(\xi)$ is an \emph{odd} continuous
function on the $2$-sphere $\mathbb{S}^2 \subset
\mathbb{R}^3$.  For any $\xi \in \mathbb{S}^2$, the
restriction of $f$ to the great circle through $\xi$ and
$-\xi$ satisfies $f(-\xi) = -f(\xi)$; by the intermediate
value theorem there exists $\xi_0$ on this great circle with
$f(\xi_0) = 0$.  Hence $\det M(\xi_0) = 0$ and the symbol
drops rank at $\xi_0$.

Since the argument is independent of the choice of~$x$,
$s_0$, and $(A_0,B_0,C_0)$, the system is never
DN-elliptic.
\end{proof}

\begin{remark}[Explicit verification]
\label{rem:n3-explicit}
For the special case $p = e_1$ and $H = \mathrm{diag}(\lambda_1,
\lambda_2, \lambda_3)$, the off-diagonal block $D$ of $M(\xi)$
vanishes and a Laplace expansion gives
$\det M(\xi) = -\frac{1}{4}(\lambda_2 - \lambda_3)
\xi_1\xi_2\xi_3$,
which vanishes on the three coordinate great circles.
For general $(p,H)$ the determinant is a nontrivial
homogeneous cubic whose zero locus is a real algebraic curve in
$\mathbb{S}^2$, but the parity constraint ensures it is
\emph{always} nonempty.
\end{remark}

\begin{theorem}[No tame right inverse in dimension $3$]
\label{thm:n3-no-tame}
Let $(M^3,h)$ be a closed 3-manifold.  For generic background
data $(s_0,A_0,B_0,C_0)$ with $C_0 \neq 0$, the linearised
extended seam map
\[
  L := DF_{(s_0,\id,A_0,B_0,C_0)}
  \colon H^{s+2} \oplus H^{s+1}(TM) \oplus (H^s)^3
  \;\longrightarrow\;
  H^s(\Sym^2 T^*M)
\]
does not admit a right inverse bounded in Sobolev norms.
In particular, no Nash--Moser iteration with finite
derivative loss can establish smooth local surjectivity
of the nonlinear map.
\end{theorem}

\begin{proof}
The characteristic variety of the linearised system is
\[
  \Sigma := \{(x,\xi) \in S^*M : \det M(\xi) = 0\}.
\]
By Theorem~\ref{thm:n3-parity}, $\det M(\xi)$ is
a homogeneous (odd) cubic in~$\xi$;
for generic data its zero locus on each fibre $S^2_x$
is a smooth algebraic curve, giving $\Sigma$
codimension~$1$ in the $5$-dimensional cosphere
bundle~$S^*M$.  At a smooth point of $\Sigma$, the
rank of~$M(\xi)$ drops by exactly~$1$ (generically),
so the system has \emph{simple characteristics of
real principal type}.

Consider the formal adjoint $L^*$, which is an
overdetermined DN system with symbol $M(\xi)^T$.
At $(x_0,\xi_0) \in \Sigma$ there exists
$v \neq 0 \in \Sym^2(\mathbb{R}^3)$ with
$M(\xi_0)^T v = 0$.  By H\"ormander's propagation
of singularities theorem for operators of real principal
type~\cite{Hormander1985III}, the wave-front set of
any distribution $u$ satisfying $L^* u \in C^\infty$
propagates along the bicharacteristic flow within
$\Sigma$.  Since $\Sigma$ has codimension~$1$,
the bicharacteristics foliate an open dense subset
of~$\Sigma$.

In particular, for any $s \in \mathbb{R}$ there exist
non-smooth solutions $u \in H^s_{\mathrm{loc}} \setminus
H^{s+\varepsilon}$ to $L^* u \in C^\infty$.
If $L$ admitted a bounded right inverse
$R \colon H^s \to H^{s+\ell}$ for some finite~$\ell$,
then $L^*$ would admit a bounded left inverse
$R^* \colon H^{-s-\ell} \to H^{-s}$, giving the
a~priori estimate
$\|u\|_{H^{-s}} \le C\,\|L^* u\|_{H^{-s-\ell}}$
for all $u$.  But the propagating singularities violate
this estimate: one can construct $u$ with
$\|u\|_{H^{-s}} = \infty$ yet
$L^* u \in C^\infty \subset H^{-s-\ell}$, a
contradiction.
\end{proof}

\begin{remark}[Parity threshold and jet order $3$]
\label{rem:parity-threshold}
The obstruction in Theorems~\ref{thm:n3-parity}
and~\ref{thm:n3-no-tame} is fundamentally a \emph{parity
phenomenon}: with $3$ diffeomorphism columns (each linear
in~$\xi$) and $3$ generator columns (constant in~$\xi$),
the symbol determinant is a homogeneous polynomial of
odd degree $3$ in $\xi \in \mathbb{R}^3$, forcing a zero.

Adding a \emph{fourth} free generator column (e.g.\ from
a third-order equivariant such as
$D\cdot(\Delta s\,ds{\otimes}ds)$) changes the symbol matrix
from $6{\times}6$ to $6{\times}7$.  The maximal minors
obtained by dropping a \emph{diffeomorphism} column
then have only $2$ linear and $4$ constant columns,
yielding a determinant of even degree~$2$ in~$\xi$.
A nonzero homogeneous polynomial of even degree on
$\mathbb{R}^3$ \emph{can} be sign-definite on
$\mathbb{S}^2$, so the parity obstruction disappears.

Thus, the passage from jet order $k = 2$
(three generators) to $k = 3$ (four or more generators)
is the precise threshold at which the symbol
parity obstruction is broken.  Whether the additional
generators at order~$3$ actually yield local universality
in dimension~$3$ remains an open and delicate question,
but the \emph{topological} barrier is lifted.
\end{remark}

\begin{remark}[Dimension heuristic at order $3$ in dimension $3$]
\label{rem:dim-heuristic-order3}
For $n = 3$ at jet order $k = 3$, the source of the extended seam map
consists of the seam ($1$ function), a diffeomorphism ($3$ functions),
and the five coefficient functions $\alpha,\beta,\gamma,\delta_1,\delta_2$
($5$ functions), giving $1 + 3 + 5 = 9$ scalar functions of $3$ variables.
The target $\Sym^2_+(T^*M)$ has $\frac{3 \cdot 4}{2} = 6$ independent
components.  Thus the system is formally \emph{overdetermined}
($9 > 6$), and no algebraic jet-dimension deficit obstructs
surjectivity.  After absorbing the seam variation into the
diffeomorphism columns (as in Theorem~\ref{thm:n3-parity}), the
effective system is $8 \times 6$, still overdetermined.  The open
question is therefore purely analytic: whether the extended
Douglis--Nirenberg symbol at order~$3$ is generically surjective.
\end{remark}

\begin{remark}[A symbol-level criterion at order $3$ in dimension $3$]
\label{rem:order3-symbol-criterion}
At a point $x\in M$ fix a background seam $s_0$ and denote
$p:=\nabla s_0(x)$, $H:=\nabla^2 s_0(x)$, and
$T_i:=\Pi_{\nabla^3,i}(j^3 s_0)(x)\in\Sym^2(T_x^*M)$.
After absorbing the principal symbol of the seam-variation into the
diffeomorphism columns as in the proof of
Theorem~\ref{thm:n3-parity}, the principal symbol of the order-$3$
extended system at $x$ has the form
\[
M(\xi)(\tilde X,\dot\alpha,\dot\beta,\dot\gamma,\dot\delta_1,\dot\delta_2)
= \xi\odot \tilde X
  + \dot\alpha\,h
  + \dot\beta\,(p\otimes p)
  + \dot\gamma\,H
  + \dot\delta_1\,T_1
  + \dot\delta_2\,T_2,
\qquad \xi\in T_x^*M\setminus\{0\}.
\]
Set $V:=\mathrm{span}\{h,\,p\otimes p,\,H,\,T_1,\,T_2\}\subset\Sym^2(T_x^*M)$.
If $\dim V=5$, then $V^\perp$ is one-dimensional (with respect to the
$h$-induced inner product on $\Sym^2$); fix $0\neq v\in V^\perp$.
Then for $\xi\neq 0$ the map $M(\xi)$ is surjective onto
$\Sym^2(T_x^*M)$ if and only if $v\,\xi\neq 0$ (viewing $v$ as a symmetric
endomorphism via $h$).  In particular, if $\det(v)\neq 0$ then the symbol
is surjective for \emph{all} $\xi\neq 0$.

Indeed, surjectivity fails precisely when the $3$-dimensional subspace
$S_\xi:=\{\xi\odot Y:Y\in T_xM\}$ lies in $V$; since $v\perp V$, this is
equivalent to $v\perp S_\xi$.  But $\langle v,\xi\odot Y\rangle=2\langle v\xi,Y\rangle$
for all $Y$, so $v\perp S_\xi$ holds if and only if $v\xi=0$.
\end{remark}

\begin{openproblem}[Universality at jet order $3$]
\label{prob:universality}
At order $\le 2$ we have the following picture:
\begin{center}
\renewcommand{\arraystretch}{1.2}
\begin{tabular}{lll}
\textbf{Setting} & \textbf{Result} & \textbf{Theorem}\\
\hline
$n=2$, genus~$0$    & universal
  & Prop.~\ref{prop:genus0}\\
$n=2$, genus~$1$    & locally universal
  & Thm.~\ref{thm:torus}\\
$n=2$, any genus, $\gamma\neq 0$ & non-elliptic (DN symbol kernel)
  & Thm.~\ref{thm:no-go-surface}\\
$n=3$, free coefficients & no tame right inverse
  & Thm.~\ref{thm:n3-no-tame}\\
$n\ge 4$, free coefficients & nowhere dense
  & Thm.~\ref{thm:free-coeff-no-go}\\
Any fixed rule, $n\ge 3$ & infinite codimension
  & Prop.~\ref{prop:jet-deficit}\\
\end{tabular}
\end{center}
For genus $g\ge 2$, local universality on surfaces remains open.
At jet order $k = 3$ the parity obstruction disappears
(Remark~\ref{rem:parity-threshold}).  The remaining
question is:
\begin{enumerate}[label=(\alph*)]
\item Do the $\OO(n)$-equivariant generators at order $3$
  yield local universality in dimension $n = 3$?
\item What is the minimal jet order $k$ for universality on
  general closed $n$-manifolds?
\end{enumerate}
\end{openproblem}

\begin{remark}[Why the no-go theorem matters]
The non-universality results above show that seam-generated
metrics form a \emph{geometrically distinguished} subclass
of all Riemannian metrics---analogous to how K\"ahler metrics,
controlled by a single plurisubharmonic potential, form a rigid
subclass within Hermitian metrics.  Far from being a weakness,
this rigidity is what makes the framework useful: the seam
representation imposes \emph{structural constraints} (the
degeneration locus, the Fredholm theory, the Gauss--Bonnet
identity) that would not hold for arbitrary metrics.
\end{remark}

Despite the global obstructions, the elliptic theory developed
in \S\ref{sec:ellipticity} yields a \emph{local} result:

\begin{proposition}[Local surjectivity of the seam map]
\label{prop:local-surj}
Let $(M^n,h)$ be compact without boundary, let $\gamma \neq 0$,
and suppose the coefficients $\alpha,\beta,\gamma$ depend smoothly on
$(s, |\nabla s|^2)$.  Fix a seam $s_0$ such that $g_0 = \Phi(s_0;h)$
is a Riemannian metric.  Then the image of~$\Phi$ contains an
$L^2$-open neighbourhood of~$g_0$ in
$\Met(M)/\ker D\Phi_{s_0}$.

More precisely: for every symmetric 2-tensor $\dot g$ in the image
of $D\Phi_{s_0}$, there exists $\varepsilon > 0$ and a smooth
family $s_t$ ($|t| < \varepsilon$) with $s_0 = s_0$ and
$\Phi(s_t;h) = g_0 + t\,\dot g + O(t^2)$.  When
$\dot g \in \image(D\Phi_{s_0})^\perp$, the deformation is
obstructed to first order.
\end{proposition}

\begin{proof}
The normal operator $N_{s_0} := (D\Phi_{s_0})^* \circ D\Phi_{s_0}
\colon H^{k+4}(M) \to H^k(M)$ is a fourth-order elliptic scalar
operator with principal symbol $\gamma^2|\xi|^4$
(Proposition~\ref{prop:fredholm}), hence Fredholm of index zero.
Since $N_{s_0}$ is self-adjoint, its image equals
$(\ker N_{s_0})^\perp$.
The implicit function theorem for Banach spaces, applied to
the map $s \mapsto \Phi(s;h)$ at $s_0$, then gives local
surjectivity onto the image of $D\Phi_{s_0}$ in a Sobolev
neighbourhood of $g_0$.
Smoothness of the family $s_t$ follows from elliptic regularity.
\end{proof}

In concrete terms: when $\gamma \neq 0$ and the kernel of the
normal operator is trivial (which is generic), every nearby metric
in the image direction can be reached by adjusting the seam.

% ================================================================
\section{Discrete Rules and Quadrature}
\label{sec:discrete}
% ================================================================

On a graph $G = (V,E)$ with background edge lengths
$\ell_0 \colon E \to \mathbb{R}_+$, a seam is a function
$s \colon V \to \mathbb{R}$.  A discrete rule assigns new edge lengths
$\ell_s(u,v)$ and the resulting geometry is the shortest-path metric
$d_s$ on the weighted graph.

\begin{proposition}[Discrete rules as quadrature of continuous rules]
\label{prop:quadrature}
Let $\Phi_\text{cont}$ be a continuous conformal-type rule
$g = f(s)\,h$ for some function $f > 0$.  The corresponding
discrete graph rule obtained by trapezoidal quadrature is
\begin{equation}
\ell_s(u,v) = \ell_0(u,v) \cdot
\frac{f(s(u)) + f(s(v))}{2}.
\end{equation}
This is the unique endpoint rule that:
\begin{enumerate}[label=(\roman*)]
\item is symmetric in $(u,v)$,
\item recovers the continuous conformal length element
  $\int_\gamma f(s)\,d\ell_0$ to second order along each edge
  (trapezoidal consistency), and
\item reduces to $\ell_0(u,v)$ when $s$ is constant with $f(s) = 1$.
\end{enumerate}
The conformal graph rule (Example~\ref{ex:conformal}) corresponds to
$f(s) = e^s$, while the gradient graph rule corresponds to
$f(s) = |\nabla s|(x)$ evaluated at vertex positions.
\end{proposition}

\begin{remark}[Discrete rules inherit the linearised structure]
The linearisation of the discrete conformal graph rule at $s = 0$
yields the cotangent Laplacian, which is
the discrete avatar of the linearised continuous conformal rule.
This is a concrete instance of the general principle that
\emph{the deformation complex of the continuous rule discretises to
the deformation complex of the discrete rule}.
\end{remark}

% ================================================================
\section{Higher-Order Rules and the Jet Hierarchy}
\label{sec:higher}
% ================================================================

\subsection{The 4-jet rule}

At order 4, the jet $j^4 s$ includes the totally symmetric tensor
$T_{ijkl} = \nabla_i \nabla_j \nabla_k \nabla_l s$.  The
$\OO(n)$-invariant \emph{scalar} formed by full double-contraction
with the inverse metric is
\begin{equation}
K_4[s] := g^{ij}g^{kl} T_{ijkl}
= \Delta^2 s + \text{(lower-order curvature corrections)}.
\end{equation}
This scalar controls the first-order correction to sublevel-set volume
in phase-space applications.

\begin{proposition}[Universality of the 4-jet correction sign]
For any non-Gaussian (non-quadratic) normalisable seam
$e^{-s} \in L^1$, the spherically averaged double-trace
$K_4[s]$ is strictly positive.
This is a consequence of the AM--GM inequality applied to even-order
jet tensors under spherical integration.
\end{proposition}

\subsection{The order-3 classification}

At order $k = 3$, the jet $j^3 s$ introduces the totally symmetric
3-tensor $\nabla^3 s \in S^3(T^*M)$.  With the $\OO(n)$-action
fixed, the new irreducible summands in $S^3(\mathbb{R}^n)$ yield
exactly two additional equivariant maps to $\Sym^2(\mathbb{R}^n)$
beyond those available at order~2.

\begin{proposition}[Order-3 generators]
\label{prop:order3}
The space of smooth $\OO(n)$-equivariant maps
$J^3(M,\mathbb{R}) \to \Sym^2(T^*M)$ that are at most linear in
the 3-jet tensor $\nabla^3 s$ is spanned by the three order-2
generators of Theorem~\ref{thm:classification} together with two new
terms:
\begin{enumerate}[label=(\roman*)]
\item $\Pi_{\nabla^3,1} \colon j^3 s \mapsto
  (\nabla_i s)(\nabla_j s)(\nabla_k s)\,g^{kl}\,(\nabla^3 s)_{ljm}
  \;+\; (i \leftrightarrow j)$,
  a cubic-in-gradient term contracted against the 3-jet,
\item $\Pi_{\nabla^3,2} \colon j^3 s \mapsto
  g^{kl}(\nabla^3 s)_{ikl}\,(\nabla_j s)
  \;+\; (i \leftrightarrow j)
  = (\nabla_i \Delta s)(\nabla_j s) + (i \leftrightarrow j)$,
  the symmetrised product of the gradient of the Laplacian and the
  gradient.
\end{enumerate}
The resulting \emph{five-parameter family}
\[
\Phi^{(3)}
= \alpha\, h
  + \beta\, ds \otimes ds
  + \gamma\, \nabla^2 s
  + \delta_1\, \Pi_{\nabla^3,1}
  + \delta_2\, \Pi_{\nabla^3,2}
\]
(with coefficient functions of the $\OO(n)$-invariant scalars)
classifies all isotropic natural metric rules of order~$\le 3$
that are at most linear in the 3-jet.
\end{proposition}

\begin{proof}[Proof sketch]
The symmetric cube $S^3(\mathbb{R}^n)$ decomposes under $\OO(n)$
into $\mathbb{R}^n \oplus W$, where $W$ is the trace-free part
(the harmonic cubic tensors).  New equivariant maps
$J^3 \to \Sym^2$ that are linear in the $3$-jet must factor through
$\Hom_{\OO(n)}(S^3(\mathbb{R}^n) \otimes S^1(\mathbb{R}^n),\;
\Sym^2(\mathbb{R}^n))$,
since the gradient $\nabla s \in S^1$ is the only available
vector at lower jet order.  Under $\OO(n)$ we have the
tensor product decomposition
\[
S^3 \otimes S^1
\;=\;
(\mathbb{R}^n \oplus W) \otimes \mathbb{R}^n
\;\cong\;
S^2 \oplus \Lambda^2 \oplus (W \otimes \mathbb{R}^n).
\]
The multiplicity of the $\Sym^2$ irreducible in this decomposition
is exactly~$2$ (one copy from the $S^2$ summand in
$\mathbb{R}^n \otimes \mathbb{R}^n$, one from the trace contraction
$W \to \mathbb{R}^n$ followed by symmetrisation with the gradient).
These correspond to $\Pi_{\nabla^3,1}$ and $\Pi_{\nabla^3,2}$
respectively.  No further linearly independent equivariants exist.
\end{proof}

\subsection{The filtration of rules by jet order}

Rules naturally organise into a filtration:
\[
\underbrace{\text{order } 0}_{\text{pointwise}}
\;\subset\;
\underbrace{\text{order } 1}_{\text{gradient}}
\;\subset\;
\underbrace{\text{order } 2}_{\text{Hessian, 3 generators}}
\;\subset\;
\underbrace{\text{order } 3}_{\text{5 generators}}
\;\subset\;
\underbrace{\text{order } 4}_{\text{biharmonic}}
\;\subset\; \cdots
\]
At each order, the classification generalises: the space of isotropic
equivariant maps is determined by the representation theory of
$\OO(n)$ on symmetric tensor products, and is computable
(e.g.~via LiE, Magma, or symbolic tensor algebra packages).

\begin{openproblem}[Complete classification at order $k$]
For arbitrary $k$, classify all smooth $\OO(n)$-equivariant natural
operators $J^k(M,\mathbb{R}) \to \Sym^2(T^*M)$.  Determine the
dimension of the space of generators and its growth rate as a function
of $k$ and $n$.  Proposition~\ref{prop:order3} gives $\dim = 5$ at
$k = 3$; the pattern at higher orders is open.
\end{openproblem}

% ================================================================
\section{Seam Geometry as a Unifying Language}
\label{sec:unifying}
% ================================================================

We collect the main identifications that connect the jet-theoretic
framework to the motivating applications.

\medskip
\begin{center}
\renewcommand{\arraystretch}{1.3}
\begin{tabular}{llll}
\hline
\textbf{Application} & \textbf{Seam} & \textbf{Rule (type)}
& \textbf{Degeneration locus} \\
\hline
Conformal geometry
  & $s \in C^\infty(M)$
  & Conformal ($\alpha = e^{2s}$)
  & --- (none) \\
Score-based diffusion
  & $\log p(x,t)$
  & Gradient ($\alpha = |\nabla s|^2$)
  & Modes/saddles of $p$ \\
Information geometry
  & Log-partition $A(\theta)$
  & Hessian ($\gamma = 1$)
  & Degenerate models \\
Optimal transport
  & Brenier potential $\varphi$
  & Hessian ($\gamma = 1$)
  & Singular support \\
Phase-space uncertainty
  & $\sum x_i^2/\sigma_i^2$
  & Hessian + 4-jet ($\gamma = 1$)
  & --- (globally nondeg.) \\
\hline
\end{tabular}
\end{center}
\medskip

In every case, the seam is a single scalar field, the rule is a
natural differential operator captured by the Universal Decomposition
(or its complex/discrete analogue), and the geometry of interest
concentrates near the degeneration locus.

% ================================================================
\section{Open Problems and Future Directions}
\label{sec:open}
% ================================================================

\begin{openproblem}[Universality at jet order $3$]
The universality landscape at order~$\le 2$ is fully resolved
(see the summary table in Open Problem~\ref{prob:universality}).
The parity obstruction in dimension $n = 3$ disappears at
jet order~$3$ (Remark~\ref{rem:parity-threshold}),
where Proposition~\ref{prop:order3} provides five generators.
Do these generators yield local universality on
closed $3$-manifolds?  More generally, what is the minimal
jet order for universality on $n$-manifolds with $n \ge 4$?
\end{openproblem}

\begin{openproblem}[Spectral theory of the linearised rule]
For a given rule $\Phi$ (e.g., the conformal graph rule), study the
spectrum of the linearised operator $D\Phi_s$ as $s$ varies.
In the overdetermined-elliptic regime (e.g. when $\gamma(s)\neq 0$ and
the coefficients depend only on $(s,|\nabla s|^2)$), the relevant
Fredholm theory is governed by the fourth-order scalar normal operator
$N_s:=(D\Phi_s)^*D\Phi_s$ (Corollary~\ref{cor:elliptic} and
Proposition~\ref{prop:fredholm}).  Beyond this regime, characterise the
seams $s$ for which $D\Phi_s$ (or $N_s$) is Fredholm, and relate the
spectral gap of $N_s$ to the conditioning of the inverse-design
problem.
\end{openproblem}

\begin{remark}[Normal operators for basis rules]
Working with the $L^2$-pairings induced by the background metric $h$
(so that $(D\Phi_s)^*$ denotes the formal adjoint with respect to
$dV_h$), the normal operator $N_s := (D\Phi_s)^*D\Phi_s$ can be
computed explicitly for the simplest basis rules.

\smallskip\noindent
\emph{(i) Conformal rule.}
For $\Phi_{\mathrm{conf}}(s;h)=e^{2s}h$ we have
$D\Phi_{s_0}(\dot s)=2e^{2s_0}\dot s\,h$, hence
$(D\Phi_{s_0})^*T = 2e^{2s_0}\,\tr_h T$ and
\[
N_{s_0}\dot s
= (D\Phi_{s_0})^*\bigl(2e^{2s_0}\dot s\,h\bigr)
= 4n\,e^{4s_0}\dot s.
\]
Thus $N_{s_0}$ is an invertible multiplication operator, and the
conditioning of the inverse-design normal equation is controlled by
the oscillation of the seam through
$\kappa(N_{s_0}) = e^{4\,\mathrm{osc}(s_0)}$.

\smallskip\noindent
\emph{(ii) Hessian rule.}
For $\Phi_{\mathrm{Hess}}(s)=\nabla^2 s$ the linearisation is
$D\Phi(\dot s)=\nabla^2\dot s$ (independent of $s_0$).  The formal
adjoint is the double divergence $(\nabla^2)^*T = \nabla^j\nabla^i T_{ij}$,
so the normal operator is
\[
N\,f
= \nabla^j\nabla^i\nabla_i\nabla_j f
= \Delta_h^2 f
  + \Ric^{ij}(h)\,\nabla_i\nabla_j f
  + (\nabla^j\Ric_{ji}(h))\,\nabla^i f.
\]
In particular, $N$ depends only on the background geometry $h$.
On Ricci-flat backgrounds (and hence on flat tori) this simplifies to
$N=\Delta_h^2$, so the spectrum on $(\ker N)^\perp=(\mathrm{constants})^\perp$
is $\{\lambda_k(\Delta_h)^2\}$.

On the flat torus $\mathbb{T}^n=\mathbb{R}^n/\mathbb{Z}^n$ one may take
the Fourier modes $f_k(x)=e^{2\pi i k\cdot x}$ ($k\in\mathbb{Z}^n$), for which
$\Delta f_k=-(2\pi|k|)^2 f_k$ and hence
$N f_k = (2\pi|k|)^4 f_k$.  The first positive eigenvalue of $N$ is
$(2\pi)^4$, of real multiplicity $2n$.

Moreover, $N=(\nabla^2)^*\nabla^2$ is nonnegative and satisfies the
Bochner energy identity
\[
\langle f, Nf\rangle_{L^2}
= \|\nabla^2 f\|_{L^2}^2
= \|\Delta_h f\|_{L^2}^2 - \int_M \Ric(\nabla f,\nabla f)\,dV_h,
\]
which shows that $\ker N = \{f : \nabla^2 f = 0\}$ consists of functions
with parallel gradient.  In particular, if $(M,h)$ admits no
nontrivial parallel vector fields then $\ker N = \mathbb{R}$.
If $\Ric(h)\le 0$, then $\|\nabla^2 f\|_{L^2}^2\ge \|\Delta_h f\|_{L^2}^2$,
so the spectral gap of $N$ on $(\ker N)^\perp$ is bounded below by
$\lambda_1(\Delta_h)^2$.
More generally, if $\Ric(h)\ge -K h$ and $\lambda_1(\Delta_h)>K$ then
for $f\perp\ker N$ one obtains the quantitative lower bound
$\langle f,Nf\rangle_{L^2}\ge (\lambda_1(\Delta_h)-K)\,\lambda_1(\Delta_h)\,\|f\|_{L^2}^2$.
\end{remark}

\begin{openproblem}[Optimal rules for inverse design]
Given a target class of geometries $\mathcal{G}_\text{target}$ and a
background $(M,h)$, what is the ``best'' rule $\Phi$ (i.e., choice of
$\alpha, \beta, \gamma$) for inverse metric design?  Possible
optimality criteria include: minimal condition number of $D\Phi_s$,
maximal $\image(D\Phi_s)$, or convexity of the inverse-design
objective in the seam variable.
\end{openproblem}

\begin{remark}[Convexity for the pure Hessian inverse problem]
For the pure Hessian rule $\Phi(s)=\nabla^2 s$, the $L^2(h)$ least-squares
inverse-design objective
\[
\mathcal{L}(s):=\|\nabla^2 s - g_{\mathrm{target}}\|_{L^2(h)}^2
\]
is a convex quadratic functional of $s$.  Its second variation is
\[
\delta^2\mathcal{L}(s)[\dot s,\dot s] = 2\,\langle \dot s,\,N\dot s\rangle_{L^2(h)},
\qquad N=(\nabla^2)^*\nabla^2,
\]
so $\mathcal{L}$ is strictly convex on $(\ker N)^\perp$.
In particular, for backgrounds with no nontrivial parallel vector
fields (so $\ker N=\mathbb{R}$), the minimiser is unique up to adding a
constant.
\end{remark}

\begin{proposition}[Global minimiser for pure Hessian inverse design]
\label{prop:hess-global-min}
Let $(M,h)$ be compact without boundary and let
$g_{\mathrm{target}}\in L^2(M;\Sym^2T^*M)$.  Consider the functional
$\mathcal{L}(s)=\|\nabla^2 s-g_{\mathrm{target}}\|_{L^2(h)}^2$ on $H^2(M)$.
Then there exists a minimiser $s_*\in H^2(M)$ with $s_*\perp\ker N$.
It is unique in $(\ker N)^\perp$ and is characterised by the normal
equation
\[
N s_*=(\nabla^2)^*g_{\mathrm{target}}
\]
in the weak sense.  Moreover, if $g_{\mathrm{target}}$ is smooth then
$s_*$ is smooth.
\end{proposition}

\begin{proof}
Write $A:=\nabla^2:H^2(M)\to L^2(M;\Sym^2T^*M)$, so that
$\mathcal{L}(s)=\|As-g_{\mathrm{target}}\|_{L^2}^2$ and $N=A^*A$.
Since $N$ is a self-adjoint elliptic operator of order $4$,
its spectrum on $L^2$ is discrete with finite-dimensional kernel.
Restricting to $(\ker N)^\perp$, the operator $N$ is invertible and
bounded below by the first positive eigenvalue, hence
there exists $c>0$ such that
$\langle s,Ns\rangle_{L^2}\ge c\,\|s\|_{H^2}^2$ for all $s\perp\ker N$.
Therefore $\mathcal{L}$ is coercive and strictly convex on
$(\ker N)^\perp$, so it admits a unique minimiser $s_*$ there.
The Euler--Lagrange equation is the normal equation $A^*(As_*-g_{\mathrm{target}})=0$,
equivalently $Ns_*=(\nabla^2)^*g_{\mathrm{target}}$.
If $g_{\mathrm{target}}$ is smooth, elliptic regularity implies $s_*$ is smooth.
\end{proof}

\begin{remark}[Computational realisation of inverse design]
\label{rem:inverse-algorithm}
The analytical machinery developed in this paper yields two natural
algorithms for the inverse metric-design problem
$g_{\mathrm{target}} \approx \Phi(s;h)$:
\begin{enumerate}[label=(\alph*)]
\item \textbf{Linearised (Newton-type).}
  Given an initial seam $s_0$, solve the normal equation
  $(D\Phi_{s_0})^*(D\Phi_{s_0})\,\delta s
  = (D\Phi_{s_0})^*(g_{\mathrm{target}} - \Phi(s_0;h))$
  for the correction $\delta s$, then update $s_1 = s_0 + \delta s$
  and iterate.  The normal operator is fourth-order elliptic when
  $\gamma \neq 0$ (Proposition~\ref{prop:fredholm}), so
  each step reduces to a standard sparse linear solve
  (conjugate gradient on a positive-definite scalar operator).
  Local convergence is guaranteed by
  Proposition~\ref{prop:local-surj}.
\item \textbf{Nonlinear (gradient descent).}
  Minimise the loss functional
  $\mathcal{L}(s) = \|\Phi(s;h) - g_{\mathrm{target}}\|_{L^2}^2
  + \lambda\,\mu(\mathcal{D}(\Phi,s))$,
  where the second term penalises the measure of the degeneration
  locus.  The gradient $\nabla_s \mathcal{L}$ is computed via the
  chain rule through $D\Phi_s$, which is explicitly available
  from Proposition~\ref{prop:lin}.  The degeneration-locus
  regulariser prevents the optimiser from entering the non-elliptic
  regime ($\gamma \to 0$) where local surjectivity fails.
\end{enumerate}
Both approaches apply equally on continuous manifolds (with finite
element or spectral discretisation) and on graphs (with the discrete
rules of \S\ref{sec:discrete}).
\end{remark}

\begin{openproblem}[Non-isotropic rules and anisotropic geometry]
The classification in~\S\ref{sec:classification} assumes isotropy
(full $\OO(n)$-equivariance).  Relaxing to a subgroup (e.g.,
$\OO(d) \times \OO(d)$ for phase-space geometries, or $U(n)$ for
K\"ahler geometries) enlarges the space of natural operators.
Classify the resulting ``anisotropic rules'' and determine which new
geometric structures become seam-accessible.
\end{openproblem}

\begin{openproblem}[Seam selection principles]
\label{prob:seam-selection}
In applications where the geometry arises naturally (information
geometry, diffusion models), the seam is given by the problem.
In inverse-design settings, the seam must be chosen and the choice
is consequential.  Develop
\emph{seam selection criteria}: given a target degeneration locus
$Z \subset M$ and a symmetry group $G \le \operatorname{Isom}(M,h)$,
characterise the space of seams $s$ such that
\begin{enumerate}[label=(\alph*)]
\item $\mathcal{D}(\Phi, s) = Z$ for a specified rule~$\Phi$
  (\emph{faithfulness}),
\item $\varphi^* s = s$ for all $\varphi \in G$
  (\emph{symmetry inheritance}),
\item $s$ is as regular as possible, pushing singularities into
  the geometry rather than the seam (\emph{maximal smoothness}).
\end{enumerate}
Determining whether a canonical choice exists---or whether the
residual freedom is itself meaningful---is an open question.
\end{openproblem}

\begin{remark}[Two conservative directions]
The following points are not used elsewhere in the paper, but may be
useful guides for future work.
\begin{enumerate}[label=(\roman*)]
\item \emph{Affine structure for the Hessian rule.}
  For the pure Hessian rule $\Phi(s)=\nabla^2 s$ the map $s\mapsto\Phi(s)$
  is linear, so straight-line interpolation $s_t=(1-t)s_0+ts_1$ gives
  $\Phi(s_t)=(1-t)\Phi(s_0)+t\Phi(s_1)$.  It would be interesting to
  understand when this affine structure yields meaningful convexity
  statements on the \emph{metric} side (e.g., under the positivity
  constraint $\nabla^2 s_t\in\Sym^2_+$), and how it interacts with the
  optimal-transport interpretation of Hessian metrics.
\item \emph{Degenerate limits of order-$1$ rules.}
  Order-$1$ metrics $g=\alpha h+\beta\,ds\otimes ds$ can approach rank-deficient
  tensors when $\alpha\to 0$ or when the ellipticity condition
  $\alpha+\beta|\nabla s|_h^2>0$ degenerates.  Given the level-set
  foliation rigidity of Corollary~\ref{cor:levelset-foliation}, a natural
  question is whether these degenerations admit a well-posed limiting
  metric structure (collapse of level sets, or a sub-Riemannian/Carnot--Carath\'eodory
  geometry on a canonical distribution) in an appropriate convergence
  sense.
\end{enumerate}
\end{remark}

\begin{openproblem}[Seam flows and geometric evolution]
A time-dependent seam $s(x,t)$ induces a one-parameter family of
geometries $g(t) = \Phi(s(\cdot, t);\, h)$.  When does the evolution
$\partial_t s = F(s, \nabla s, \nabla^2 s)$ induce a geometrically
natural flow on the space of metrics (e.g., Ricci flow, Yamabe flow,
mean curvature flow)?  Classify the scalar PDEs whose seam-induced
metric evolution coincides with standard geometric flows.
\end{openproblem}

\begin{remark}[Conformal flows reduce to scalar seam equations]
For the conformal rule $g=e^{2s}h$ we have $\partial_t g = 2(\partial_t s)\,g$.
Hence any metric flow of the form $\partial_t g = q(g)\,g$ (a conformal
deformation) corresponds to the seam equation $\partial_t s = \tfrac12 q(g)$.
In particular, Yamabe flow $\partial_t g = -R(g)\,g$ corresponds to
$\partial_t s = -\tfrac12 R(e^{2s}h)$, i.e.
\[
\partial_t s
= -\tfrac12 e^{-2s}\bigl[R(h) - 2(n-1)\Delta_h s - (n-1)(n-2)|\nabla s|_h^2\bigr].
\]
\end{remark}

\begin{openproblem}[Machine-learned rules and jet-equivariant architectures]
In applications where the target geometry is known empirically
(e.g., from data), can one \emph{learn} the coefficient functions
$(\alpha, \beta, \gamma)$ from training data?  The Universal
Decomposition provides a strict architectural constraint:
\emph{any model constrained to represent a smooth isotropic natural
second-order local rule that is at most linear in $\nabla^2 s$ must
factor through the three basis tensors}
$h$, $ds \otimes ds$, $\nabla^2 s$
(Corollary~\ref{cor:forced-arch}).  This is not an optional
design choice but a mathematical necessity: the only freedom
lies in the coefficient functions $\alpha, \beta, \gamma$
evaluated on the invariant scalars~$\mathcal{I}$.

A concrete architecture is immediate: build a neural network whose
output layer is exactly the triple $(\alpha, \beta, \gamma)$
evaluated on $\mathcal{I}(j^k s)$.  Given a target metric
$g_{\mathrm{target}}$, one trains by back-propagating through the
linearised operator $D\Phi_s$ (which is explicitly computable and
differentiable by Proposition~\ref{prop:lin}).  The built-in
naturality bias guarantees that every learned rule remains a valid
natural operator---no spurious coordinate dependence.  Regularisation
can be imposed via the condition number of $D\Phi_s$ or the measure
of the degeneration locus (Definition~\ref{def:degen}).
\end{openproblem}

% ================================================================
\section{Conclusion}
% ================================================================

The seam--rule paradigm rests on a simple idea: geometry can be
generated from scalar fields through local, natural transformation
rules.  By recognising these rules as natural differential operators
on jet bundles, we obtain:

\begin{enumerate}
\item A \textbf{classification}: all isotropic metric rules of order
  $\le 2$ decompose into conformal, gradient-tensor, and Hessian
  components (Theorem~\ref{thm:classification}).
\item A \textbf{deformation theory}: the linearised rule is a
  differential operator whose Fredholm properties govern inverse
  design, stability, and gauge freedom
  (\S\ref{sec:linearisation}).
\item A \textbf{degeneration locus}: a universal geometric invariant
  that unifies critical points, zeros, nodes, and minima across
  diverse applications (\S\ref{sec:degeneration}).
\item A precise \textbf{universality question}: when does scalar data
  suffice to reach all of Riemannian geometry?
  (\S\ref{sec:universality}).
\end{enumerate}

The framework is deliberately minimalist: a single scalar field, a
natural operator, and a background structure.  Yet it encompasses
conformal geometry, Hessian geometry, information geometry, optimal
transport, and diffusion models---all as
instances of one universal decomposition.

The central open problem is universality: whether every Riemannian
metric on a compact manifold can be reached (up to diffeomorphism) by
some seam under the full three-term family of order-$2$ natural
operators.  An affirmative answer would show that Riemannian
geometry admits a scalar-generating representation up to
diffeomorphism---a structural reduction analogous to the
potential-theoretic description of K\"ahler metrics.

\backmatter

\bmhead{Acknowledgements}

Not applicable.

\section*{Declarations}

\bmhead{Funding}

Not applicable.

\bmhead{Competing interests}

The author declares no competing interests.

\bmhead{Use of large language models (LLMs)}

During preparation of this manuscript, the author used large language
models (Claude Opus~4.6, GPT-5.2, Grok~4.20, and Gemini~3.1 Pro) to
assist with drafting and editing text and with exploring alternative
expositions. The author reviewed and edited the resulting content and
takes full responsibility for the manuscript.

\bmhead{Ethics approval and consent to participate}

Not applicable.

\bmhead{Consent for publication}

Not applicable.

\bmhead{Data availability}

Not applicable.

\bmhead{Materials availability}

Not applicable.

\bmhead{Code availability}

Not applicable.

\bmhead{Author contribution}

The author wrote the manuscript.

\begin{thebibliography}{99}

\bibitem{KMS1993}
I.~Kol\'a\v{r}, P.~W.~Michor, and J.~Slov\'ak,
\textit{Natural Operations in Differential Geometry},
Springer, 1993.

\bibitem{Peetre1960}
J.~Peetre,
``Rectification \`a l'article `Une caract\'erisation abstraite des
op\'erateurs diff\'erentiels'\,''
\textit{Math.\ Scand.} \textbf{8}, 116--120 (1960).

\bibitem{Slovak1988}
J.~Slov\'ak,
``Peetre theorem for nonlinear operators,''
\textit{Ann.\ Global Anal.\ Geom.} \textbf{6}(3), 273--283 (1988).

\bibitem{Weyl1946}
H.~Weyl,
\textit{The Classical Groups: Their Invariants and Representations},
Princeton University Press, 1946.

\bibitem{Procesi2007}
C.~Procesi,
\textit{Lie Groups: An Approach through Invariants and
Representations},
Springer, 2007.

\bibitem{Hormander1985III}
L.~H\"ormander,
\textit{The Analysis of Linear Partial Differential Operators~III:
Pseudo-Differential Operators},
Springer, 1985.

\bibitem{KazdanWarner1975}
J.~L.~Kazdan and F.~W.~Warner,
``Scalar curvature and conformal deformation of Riemannian structure,''
\textit{J.\ Differential Geom.} \textbf{10}(1), 113--134 (1975).

\bibitem{Villani2009}
C.~Villani,
\textit{Optimal Transport: Old and New},
Grundlehren der mathematischen Wissenschaften~\textbf{338},
Springer, 2009.

\bibitem{Caffarelli1992}
L.~A.~Caffarelli,
``The regularity of mappings with a convex potential,''
\textit{J.\ Amer.\ Math.\ Soc.} \textbf{5}(1), 99--104 (1992).

\bibitem{Amari2016}
S.~Amari,
\textit{Information Geometry and Its Applications},
Applied Mathematical Sciences~\textbf{194},
Springer, 2016.

\bibitem{SongErmon2019}
Y.~Song and S.~Ermon,
``Generative modeling by estimating gradients of the data
distribution,''
in \textit{Advances in Neural Information Processing
Systems~32} (NeurIPS), 2019.

\bibitem{Horndeski1974}
G.~W.~Horndeski,
``Second-order scalar-tensor field equations in a
four-dimensional space,''
\textit{Int.\ J.\ Theor.\ Phys.} \textbf{10}(6), 363--384 (1974).

\end{thebibliography}

\end{document}